\documentclass[11pt]{amsart}
\usepackage[margin=1.1in]{geometry}
\usepackage[T1]{fontenc}
\usepackage{lmodern}
\usepackage{amsmath,amssymb,amsthm,mathtools}
\usepackage{xcolor}
\usepackage[colorlinks=true,linkcolor=blue!50!black,citecolor=green!40!black,urlcolor=blue!50!black]{hyperref}
\usepackage[capitalise,noabbrev]{cleveref}
\usepackage[backend=bibtex,style=alphabetic,maxbibnames=10]{biblatex}
\addbibresource{longversion.bib}

\newtheorem{theorem}{Theorem}[section]
\newtheorem{proposition}[theorem]{Proposition}
\newtheorem{lemma}[theorem]{Lemma}
\newtheorem{corollary}[theorem]{Corollary}
\newtheorem{conjecture}[theorem]{Conjecture}
\theoremstyle{definition}
\newtheorem{definition}[theorem]{Definition}
\newtheorem{example}[theorem]{Example}
\theoremstyle{remark}
\newtheorem{remark}[theorem]{Remark}
\numberwithin{equation}{section}

%% Visible work-in-progress marker; every occurrence must be resolved before posting.
\newcommand{\wip}[1]{{\color{red}\textbf{[WIP:} #1\textbf{]}}}

\newcommand{\lin}{\operatorname{lin}_e}
\newcommand{\Lead}{\operatorname{Lead}}
\newcommand{\val}{v}
\newcommand{\Lam}{\Lambda}
\newcommand{\CT}{\operatorname{CT}}
\newcommand{\Res}{\operatorname*{Res}}
\newcommand{\qbin}[2]{\genfrac{[}{]}{0pt}{}{#1}{#2}_t}
\newcommand{\hx}{\hat x}

\title[Hikita's $\star$-product: Pieri rules, dominance and valuations]{Hikita's $\star$-product on symmetric functions:\\ Pieri rules, dominance support, and the $(s-1)$-adic valuation}

\author{Rick (AI agent)}
\thanks{\wip{Authorship, affiliation, acknowledgements and the AI declaration are to be settled as for the extended abstract.}}

\date{\today}

\begin{document}

\begin{abstract}
Hikita introduced a two-parameter deformation $\star$ of the product of symmetric functions, obtained by transporting the ordinary product through the level-one polynomial representation of the affine Hecke algebra. We give an explicit rule for $e_k\star(e_{a_1}\cdots e_{a_\ell})$ for all $k$ and $\ell$, in generating-function form; its case $\ell=1$ is a Pieri rule for $e_k\star e_r$ generalising Hikita's rule for $e_1\star e_r$. Writing $e_{\lambda_1}\star\cdots\star e_{\lambda_\ell}=\sum_\mu c_{\lambda\mu}(s,t)e_\mu$, we show that $c_{\lambda\mu}$ is supported exactly on the dominance up-set of $\lambda$ with $s$-adic valuation $n(\mu)$, and that the $(s-1)$-adic valuation of $c_{\lambda\mu}$ is $\ell(\lambda)-\kappa(\lambda,\mu)$, where $\kappa$ is the maximal number of blocks into which $\lambda$ and $\mu$ split compatibly. Leading coefficients factor over blocks into connected pieces; for a full merge the connected piece is a connected-graph cumulant. The coefficients are invariant, up to a power of $s$, under complementation in a box, and all second-order leading coefficients are given in closed form via a residue formula for Green polynomials at two-part classes.
\end{abstract}

\maketitle

\setcounter{tocdepth}{1}
\tableofcontents

%% ===================================================================
\section{Introduction}\label{sec:intro}
%% ===================================================================

Hikita~\cite{Hikita25} defined a commutative, associative product $\star$ on symmetric functions over $\mathbb{Q}(q,t)$ by transporting the ordinary product through the level-one polynomial representation of the affine Hecke algebra. We write $s:=q^{-1}$. At $s=1$, $\star$ is the ordinary product. Hikita's Pieri rule \cite[Thm.~3.12]{Hikita25} reads, in our normalisation,
\[ e_1\star e_r=(1-s)[r+1]_t\,e_{r+1}+s\,e_1e_r . \]
Di Francesco and Kedem~\cite[\S8.3]{DFK19} observed that the associated $t$-deformed graded characters lose Schur positivity at generic $t$. This paper is about the structure that survives: explicit multiplication rules, triangularity, and valuations.

\subsection*{Pieri rules}
The operator $e_k\star(\cdot)$ is a sum over $k$-subsets (\cref{thm:subset}),
\[ e_k\star F=\sum_{|A|=k}\ \prod_{i\in A,\,j\notin A}\frac{x_i-tx_j}{x_i-x_j}\ \prod_{i\in A}x_i\ \cdot F\big|_{x_i\mapsto sx_i\ (i\in A)} . \]
From it we derive a closed generating function for $e_k\star(e_{a_1}\cdots e_{a_\ell})$ for all $k,\ell$ (\cref{thm:ellcol}). Since the $e_a$ generate, this determines $e_k\star F$ for every $F$. The case $\ell=1$ is a Pieri rule (\cref{cor:pieri-ekr}):
\[ e_k\star e_r=\sum_{b=0}^{k}s^b\,F_{k-b}(t^{r-b})\,e_b\,e_{r+k-b}, \]
with explicit polynomials $F_n(w)$, and at $t=0$ the coefficients form a truncated geometric distribution (\cref{cor:pieri-t0}). In the generating function, the $\ell$ columns interact only through a pairwise cross kernel, and the proof closes with a single residue identity (\cref{lem:Z}).

\subsection*{Dominance and valuations}
Write $e^\star_\lambda:=e_{\lambda_1}\star\cdots\star e_{\lambda_\ell}=\sum_\mu c_{\lambda\mu}(s,t)\,e_\mu$. A termwise degree count on the subset formula shows that $c_{\lambda\mu}\in\mathbb{Q}[s,t]$ is supported exactly on the dominance up-set of $\lambda$, with $s$-adic valuation exactly $n(\mu)$ (\cref{thm:DS}). In the rescaled basis $b_\mu=s^{n(\mu)}e_\mu$, the operator $e_k\star(\cdot)$ reduces modulo $s$ to Hall--Littlewood multiplication by $t^{-\binom k2}e_k$, so the $s$-leading coefficients form the classical $e\to$Hall--Littlewood transition matrix (\cref{thm:H,cor:dmatrix}).

At $s=1$, $\star$ is the ordinary product, and the variable $s$ enters the subset formula only through an Euler operator; so the $(s-1)^p$ Taylor piece of $e_k\star(\cdot)$ is a differential operator of order $p$. The first-order term is a symmetric biderivation with an explicit matrix (\cref{thm:bider,prop:M}), and an order-$p$ piece can glue at most $p$ existing blocks of parts. This gives the \emph{block valuation law}
\[ \val_{s-1}(c_{\lambda\mu})=\ell(\lambda)-\kappa(\lambda,\mu)\qquad(\mu\trianglerighteq\lambda,\ \mu\ne\lambda), \]
where $\kappa$ is the maximal number of compatible blocks (\cref{thm:blocklaw}); the lower bound holds for all $t$, and exactness is settled at $t=0$, where no cancellation occurs. The leading coefficients are sums over merge histories with explicit weights (\cref{thm:coarse,thm:W}). For a full merge the lead is $\frac{1-t^n}{\prod_i(1-t^{\lambda_i})}\sum_{H\text{ connected}}\prod_{ij\in H}(t^{\lambda_i\lambda_j}-1)$ (\cref{thm:G}), whose graph half is due to Do{\l}{\k{e}}ga~\cite{Dolega19}, and every lead is a sum of products of connected leads (\cref{thm:F,thm:blockmult}). The coefficients satisfy $c_{N^\ell-\lambda,N^\ell-\mu}=s^{N\binom\ell2-(\ell-1)|\lambda|}c_{\lambda\mu}$ (\cref{thm:box}). Finally, every leading coefficient at valuation $2$ is given in closed form (\cref{thm:v2,cor:v2}), via a constant-term formula and a residue evaluation of Green polynomials at two-part classes (\cref{thm:CT,thm:2pt}).

\subsection*{Organisation}
\Cref{sec:hikita} sets up Hikita's product and proves the subset formula. \Cref{sec:pieri} contains the Pieri package. \Cref{sec:DS} proves dominance support, and \cref{sec:edge} identifies the $s\to0$ edge. \Cref{sec:blocklaw,sec:leads} treat the expansion at $s=1$, \cref{sec:box} the box symmetry and the linear coefficient, and \cref{sec:v2} the second-order leading coefficients.

%% ===================================================================
\section{Hikita's product and the subset formula}\label{sec:hikita}
%% ===================================================================

\subsection{Notation}
For $m\ge0$ let $\Lam_m:=\mathbb{Q}(s,t)[x_1,\dots,x_m]^{S_m}$ and let $\Lam$ be the ring of symmetric functions over $\mathbb{Q}(s,t)$. We write $e_r$ for the elementary symmetric functions ($e_0=1$, $e_r=0$ for $r<0$), $e_\lambda=e_{\lambda_1}e_{\lambda_2}\cdots$, and
\[ E(z):=\prod_{i=1}^m(1+x_iz)=\sum_{r\ge0}e_rz^r,\qquad Q(y):=\frac{E(-ty)}{E(-y)}=\prod_{i=1}^m\frac{1-tx_iy}{1-x_iy}=\sum_{n\ge0}q_ny^n, \]
so that the $q_n$ are the one-row Hall--Littlewood functions of \cite[III (2.10)]{Macdonald95}. For a set $B$ of indices, $E(x_B;z):=\prod_{i\in B}(1+x_iz)$, and $\hx_i$ denotes the variables $x_1,\dots,x_m$ with $x_i$ omitted. We put
\[ a_{ij}:=\frac{x_i-tx_j}{x_i-x_j},\qquad c_A:=\prod_{i\in A,\ j\notin A}a_{ij},\qquad x_A:=\prod_{i\in A}x_i\qquad(A\subseteq[m]). \]
Further, $[n]=[n]_t=(1-t^n)/(1-t)$, $[n]_t!=[1][2]\cdots[n]$, $(a;t)_n=\prod_{i=0}^{n-1}(1-at^i)$, $n(\lambda)=\sum_i(i-1)\lambda_i$, and $\trianglerighteq$ is dominance order.

\subsection{The level-one representation and Hikita's product}
The symmetric group $S_m$ acts on $\mathbb{Q}(s,t)[x_1,\dots,x_m]$ by permuting variables; $s_i$ is the transposition of $x_i$ and $x_{i+1}$. Define operators on $\mathbb{Q}(s,t)(x_1,\dots,x_m)$ by
\begin{align*}
T_if&:=t\,s_if+(t-1)\,\frac{x_{i+1}}{x_i-x_{i+1}}\,(s_if-f)\qquad(1\le i\le m-1),\\
\pi f(x_1,\dots,x_m)&:=x_1\,f(x_2,\dots,x_m,sx_1),\\
Y_i&:=t^{m-i}\,T_{i-1}\cdots T_1\,\pi\,T_{m-1}^{-1}\cdots T_i^{-1}\qquad(1\le i\le m).
\end{align*}
The $T_i$ are Demazure--Lusztig operators; they preserve polynomials, satisfy $(T_i-t)(T_i+1)=0$ and the braid relations, and $T_if=tf$ whenever $s_if=f$. For $1\le k\le m$ put
\[ e_k(Y):=\sum_{b_1<\dots<b_k}Y_{b_1}Y_{b_2}\cdots Y_{b_k}, \]
each product taken in increasing order of indices from left to right. (The $Y_i$ commute, but the proof of \cref{thm:subset} does not use this.) Hikita's product $F\star G$ is defined by transporting the product through the map $F(Y)\mapsto F(Y)\bullet1$; by \cite[Lemma~3.1, Cor.~3.9]{Hikita25} this map is a linear isomorphism compatible with $m\to m-1$, so $\star$ is well defined on $\Lam$, and it is commutative and associative \cite{Hikita25}. We use $\star$ only through the following consequence of the definitions \cite[Def.~3.4, Lemma~3.3]{Hikita25}: for symmetric $F$,
\begin{equation}\label{eq:star-reading}
e_k\star F=t^{-\binom k2}\,e_k(Y)\,F .
\end{equation}

\subsection{Coxeter preliminaries}
We use standard facts about $S_m$ as a Coxeter group with generators $s_1,\dots,s_{m-1}$ \cite[\S\S1.10, 5.12]{Humphreys90}, \cite[\S2.4]{BjornerBrenti05}. Words act as compositions of operators, the rightmost factor first.
\begin{enumerate}
\item[(H1)] For a reduced word $w=s_{i_1}\cdots s_{i_\ell}$ the operator $T_w:=T_{i_1}\cdots T_{i_\ell}$ depends only on $w$, and $T_uT_v=T_{uv}$ whenever $\ell(uv)=\ell(u)+\ell(v)$.
\item[(H2)] For $J\subseteq\{s_1,\dots,s_{m-1}\}$ let $W_J$ be the parabolic subgroup and $W^J=\{w:\ell(ws)>\ell(w)\ \forall s\in J\}$. Every $w$ factors uniquely as $w=w^Jw_J$ with $w^J\in W^J$, $w_J\in W_J$, and $\ell(w)=\ell(w^J)+\ell(w_J)$.
\item[(H3)] For $K\subseteq J$, multiplication $W^J\times(W_J)^K\to W^K$ is a length-additive bijection.
\item[(H4)] If $T_jH=tH$ for all $s_j\in J$, then $T_vH=t^{\ell(v)}H$ for all $v\in W_J$.
\item[(H5)] $\pi T_j=T_{j+1}\pi$ for $1\le j\le m-2$.
\end{enumerate}
(H3) is standard: if $u\in W^J$, $v\in(W_J)^K$ and $s\in K$, then $\ell(uvs)=\ell(u)+\ell(vs)=\ell(uv)+1$; conversely, for $x\in W^K$ write $x=x^Jx_J$, and $x_Js<x_J$ for some $s\in K$ would give $\ell(xs)<\ell(x)$. (H4) is immediate from (H1). For (H5), $\pi s_jf=x_1f(x_2,\dots,x_{j+2},x_{j+1},\dots,sx_1)=s_{j+1}\pi f$ because $j+2\le m$, and $\pi$ turns multiplication by $x_{j+1}/(x_j-x_{j+1})$ into multiplication by $x_{j+2}/(x_{j+1}-x_{j+2})$.

Let $J_n:=\{s_i:i\neq n\}$, so $W_{J_n}=S_n\times S_{m-n}$, and the coset $wW_{J_n}$ is determined by $D=w([n])$. Its minimal element $w^D$ is increasing on $[n]$ and on $[n+1,m]$, with $\ell(w^D)=\sum_l(d_l-l)$ for $D=\{d_1<\dots<d_n\}$. The word
\[ w(D):=(s_{d_1-1}\cdots s_1)(s_{d_2-1}\cdots s_2)\cdots(s_{d_n-1}\cdots s_n) \]
sends $l\mapsto d_l$ for $l\le n$: the blocks with index $l'>l$ fix $l$, block $l$ sends $l$ to $d_l$, and the blocks $l'<l$ only move points $\le d_{l'}<d_l$. It has length $\le\sum_l(d_l-l)$, so $w(D)=w^D$ and the word is reduced. We write $T_{w(D)}$ for the corresponding operator and
\[ \sigma^{(k)}:=\sum_{|D|=k}T_{w(D)} . \]

\subsection{The parabolic kernel}

\begin{lemma}\label{lem:KL}
Let $0\le n\le m$, $1\le c<d_1<\dots<d_n\le m$, and let $H$ satisfy $T_jH=tH$ for all $j\neq n$. Then
\[ (T_cT_{c+1}\cdots T_{m-1})^{-1}\,T_{w(D)}H=t^{\,n-(m-c)}\,T_{w(D-1)}H,\qquad D-1:=\{d_1-1,\dots,d_n-1\}. \]
\end{lemma}
\begin{proof}
Let $x:=s_cs_{c+1}\cdots s_{m-1}$; it is reduced of length $m-c$, sends $m\mapsto c$ and $i\mapsto i+1$ for $c\le i\le m-1$, and fixes $i<c$. Let $y:=w(D-1)$; the word is valid since $d_l-1\ge c+l-1\ge l$, and $\ell(y)=\ell(w^D)-n$. As $m\notin D-1$, the value $m$ is the largest on the increasing block $[n+1,m]$, so $y(m)=m$. Since $x$ is increasing on $[1,m-1]$, the inversions of $y$ among positions $<m$ survive in $xy$; at position $m$, $xy(m)=c$ is inverted with exactly those $i<m$ having $y(i)\ge c$, and there are $m-c$ of them. Hence $\ell(xy)=\ell(y)+m-c$. Finally $xy([n])=x(D-1)=D$ because every $d_l-1\ge c$, so $xy=w^Dv$ with $v\in W_{J_n}$ and $\ell(v)=\ell(xy)-\ell(w^D)=m-c-n$. By (H1) and (H4), $T_xT_yH=T_{xy}H=T_{w^D}T_vH=t^{m-c-n}T_{w^D}H$. Apply $T_x^{-1}$.
\end{proof}

\begin{proposition}[Parabolic kernel]\label{prop:Ak}
For symmetric $F$ and $1\le b_1<\dots<b_k\le m$, $Y_{b_1}Y_{b_2}\cdots Y_{b_k}F=t^{\binom k2}\,T_{w(\{b_1,\dots,b_k\})}\,\pi^kF$. Hence
\[ e_k(Y)F=t^{\binom k2}\,\sigma^{(k)}\pi^kF . \]
\end{proposition}
\begin{proof}
Note that $\pi^nF=x_1\cdots x_n\,F(x_{n+1},\dots,x_m,sx_1,\dots,sx_n)$ is symmetric in $x_1,\dots,x_n$ and in $x_{n+1},\dots,x_m$, so $T_j\pi^nF=t\pi^nF$ for $j\neq n$. We induct on the number of factors, the empty product being trivial. Let $c<d_1<\dots<d_n$ and assume $Y_{d_1}\cdots Y_{d_n}F=t^{\binom n2}T_{w(D)}H$ with $H=\pi^nF$. By \cref{lem:KL},
\[ Y_c\,t^{\binom n2}T_{w(D)}H=t^{\binom n2}t^{m-c}\,T_{c-1}\cdots T_1\,\pi\,(T_c\cdots T_{m-1})^{-1}T_{w(D)}H=t^{\binom n2+n}\,T_{c-1}\cdots T_1\,\pi\,T_{w(D-1)}H. \]
All indices in the word $w(D-1)=\prod_l(s_{d_l-2}\cdots s_l)$ lie in $[1,m-2]$, so by (H5), $\pi T_{w(D-1)}=\prod_l(T_{d_l-1}\cdots T_{l+1})\,\pi$. Now $T_{c-1}\cdots T_1\prod_l(T_{d_l-1}\cdots T_{l+1})=T_{w(\{c\}\cup D)}$, $\pi H=\pi^{n+1}F$, and $\binom n2+n=\binom{n+1}2$.
\end{proof}

\subsection{The $k$-set kernel}
Let $\sigma_{[l,m]}:=\sum_{a=l}^mT_{a-1}\cdots T_l$ be the partial symmetriser on $x_l,\dots,x_m$.

\begin{lemma}\label{lem:sym1}
Let $F$ be symmetric in $x_2,\dots,x_m$, with coefficients in any field containing $\mathbb{Q}(t)$, and let $F^{(i)}$ be $F$ with $x_1$ and $x_i$ swapped. Then $\sigma_{[1,m]}F=\sum_{i=1}^mF^{(i)}\prod_{j\neq i}a_{ij}$.
\end{lemma}
\begin{proof}
Induction on $m$; $m=1$ is trivial. Since $s_1F=F^{(2)}$, a direct computation gives $T_1F=a_{21}F^{(2)}+(a_{12}-1)F^{(1)}$. Write $\sigma_{[1,m]}=1+\sigma_{[2,m]}T_1$. The function $G:=T_1F$ is symmetric in $x_3,\dots,x_m$, so by induction (over the field with $x_1$ adjoined) $\sigma_{[2,m]}G=\sum_{i\ge2}G|_{x_2\leftrightarrow x_i}\prod_{j\ge2,\,j\neq i}a_{ij}$. Under $x_2\leftrightarrow x_i$, $F^{(2)}\mapsto F^{(i)}$, $F^{(1)}$ is fixed, $a_{21}\mapsto a_{i1}$ and $a_{12}\mapsto a_{1i}$. Hence
\[ \sigma_{[1,m]}F=\sum_{i\ge2}F^{(i)}\prod_{j\neq i}a_{ij}+F^{(1)}\Bigl[1+\sum_{i\ge2}(a_{1i}-1)\prod_{j\ge2,\,j\neq i}a_{ij}\Bigr]. \]
The function $R(x):=\prod_{j\ge2}(x-tx_j)/(x-x_j)$ tends to $1$ at $\infty$ and has simple poles at $x=x_i$ with residues $(1-t)x_i\prod_{j\ge2,j\neq i}a_{ij}$; so $R(x)=1+\sum_{i\ge2}\frac{(1-t)x_i}{x-x_i}\prod_{j\ge2,j\ne i}a_{ij}$. Evaluating at $x=x_1$ and using $(1-t)x_i/(x_1-x_i)=a_{1i}-1$ shows that the bracket equals $R(x_1)=\prod_{j\ge2}a_{1j}$.
\end{proof}

\begin{proposition}[$k$-set kernel]\label{prop:Kk}
Let $G$ be symmetric in $x_1,\dots,x_k$ and in $x_{k+1},\dots,x_m$, and let $G^{(A)}$ denote $G$ with the variables $x_A$ in the first $k$ slots and $x_{[m]\setminus A}$ in the last $m-k$. Then
\[ \sigma^{(k)}G=\sum_{|A|=k}c_A\,G^{(A)} . \]
\end{proposition}
\begin{proof}
Let $K_l:=\{s_{l+1},\dots,s_{m-1}\}$. The minimal coset representatives of $W_{K_{l-1}}/W_{K_l}$ are $s_{a-1}\cdots s_l$ ($l\le a\le m$), so applying (H3) along $K_k\subset K_{k-1}\subset\dots\subset K_0$ gives $\sum_{u\in W^{K_k}}T_u=\sigma_{[1,m]}\sigma_{[2,m]}\cdots\sigma_{[k,m]}$. Applying (H3) to $K_k\subset J_k$, with $(W_{J_k})^{K_k}=S_k$, gives $\sum_{u\in W^{K_k}}T_u=\sigma^{(k)}\sum_{v\in S_k}T_v$. For $G$ as in the statement, (H4) gives $\sum_{v\in S_k}T_vG=[k]_t!\,G$, so
\begin{equation}\label{eq:chain}
\sigma_{[1,m]}\cdots\sigma_{[k,m]}G=[k]_t!\,\sigma^{(k)}G .
\end{equation}
On the other hand, if $G$ is merely symmetric in $x_{k+1},\dots,x_m$, then iterating \cref{lem:sym1} gives
\begin{equation}\label{eq:ordered}
\sigma_{[1,m]}\cdots\sigma_{[k,m]}G=\sum_{(a_1,\dots,a_k)\ \text{distinct}}G(x_{a_1},\dots,x_{a_k};\text{rest})\ \omega(a),\qquad \omega(a):=\prod_{l=1}^k\ \prod_{j\notin\{a_1,\dots,a_l\}}a_{a_lj}.
\end{equation}
Indeed, by induction on $k$ in the variables $x_2,\dots,x_m$, $H:=\sigma_{[2,m]}\cdots\sigma_{[k,m]}G$ is the analogous sum over distinct tuples in $[2,m]$; it is symmetric in $x_2,\dots,x_m$, since permuting those variables permutes the terms; and \cref{lem:sym1} applied to $H$, followed by $x_1\leftrightarrow x_{a_1}$, turns tuples in $[2,m]$ into tuples in $[m]\setminus\{a_1\}$.

Now let $G$ be symmetric in both groups. Then $G(x_{a_1},\dots,x_{a_k};\text{rest})=G^{(A)}$ with $A=\{a_1,\dots,a_k\}$, and
\begin{equation}\label{eq:omega-split}
\omega(a)=c_A\prod_{1\le l<l'\le k}a_{a_la_{l'}} .
\end{equation}
By the Poincar\'e identity $\sum_{w\in S_k}w\bigl(\prod_{i<j\le k}a_{ij}\bigr)=[k]_t!$ \cite[III (1.4)]{Macdonald95} (which is also \eqref{eq:chain}=\eqref{eq:ordered} at $m=k$, $G=1$), summing over the orderings of $A$ gives $[k]_t!\,c_AG^{(A)}$. Compare with \eqref{eq:chain}.
\end{proof}

\begin{theorem}[Subset formula]\label{thm:subset}
For every $m\ge0$, $k\ge1$ and $F\in\Lam_m$,
\[ e_k\star F=t^{-\binom k2}e_k(Y)F=E_kF:=\sum_{|A|=k}c_A\,x_A\,F\big|_{x_i\mapsto sx_i\ (i\in A)} . \]
\end{theorem}
\begin{proof}
By \cref{prop:Ak,prop:Kk} with $G=\pi^kF$, since $(\pi^kF)^{(A)}=x_A\,F|_{x_i\mapsto sx_i\,(i\in A)}$; the first equality is \eqref{eq:star-reading}. For $m<k$ both sides vanish.
\end{proof}

\begin{remark}
For $k=1$ the subset formula gives back Hikita's rule $e_1\star e_r=(1-s)[r+1]e_{r+1}+se_1e_r$ \cite[Thm.~3.12]{Hikita25}; see \cref{ex:k1}. \Cref{prop:Kk} is the level-one form of the classical expression of the Hecke symmetriser as $\sum_ww\circ\prod_{i<j}a_{ij}$, restricted to minimal coset representatives.
\end{remark}

\begin{lemma}\label{lem:Ek1}
$E_k(1)=\sum_{|A|=k}c_Ax_A=e_k$.
\end{lemma}
\begin{proof}
Grouping $w\in S_m$ by $w(\{1,\dots,k\})=A$ and applying the Poincar\'e identity to $S_k\times S_{m-k}$ shows that $\sum_{|A|=k}c_Ax_A$ is the Hall--Littlewood function $P_{(1^k)}(x;t)$ as defined in \cite[III (2.1)--(2.2)]{Macdonald95}, and $P_{(1^k)}=e_k$ \cite[III (2.8)]{Macdonald95}.
\end{proof}

\begin{lemma}[Stability]\label{lem:stable}
For $F\in\Lam_m$, $(E_kF)|_{x_m=0}=E_k(F|_{x_m=0})$, the right side computed in $m-1$ variables. Hence the $e$-coefficients of any expression built from the $E_k$ and $1$ are independent of $m$ once $m$ is at least the degree.
\end{lemma}
\begin{proof}
Terms with $m\in A$ vanish because of the factor $x_A$; for $m\notin A$, $a_{im}|_{x_m=0}=1$. Restriction $\Lam_m\to\Lam_{m-1}$ maps $e_\nu\mapsto e_\nu$ and is injective in degrees $\le m-1$.
\end{proof}

\subsection{Two consequences of the kernel}
The next two lemmas drive every computation in \cref{sec:pieri}. For a function $\varphi$ of one variable put $S^\varphi_k(x_1,\dots,x_m):=\sum_{|A|=k}c_A\prod_{a\in A}\varphi(x_a)$.

\begin{lemma}[Peeling one variable]\label{lem:peel}
For $m\ge1$ and $k\ge1$, $[k]\,S^\varphi_k(x)=\sum_{i=1}^m\varphi(x_i)\prod_{j\neq i}a_{ij}\;S^\varphi_{k-1}(\hx_i)$.
\end{lemma}
\begin{proof}
By \eqref{eq:omega-split} and the Poincar\'e identity, $[k]_t!\,S^\varphi_k=\sum_a\omega(a)\prod_l\varphi(x_{a_l})$, the sum over distinct $k$-tuples. Separate $a_1=i$: $\omega(a)=\prod_{j\neq i}a_{ij}\cdot\omega'(a_2,\dots,a_k)$, where $\omega'$ is the weight of \eqref{eq:ordered} for the variables $\hx_i$. Summing over $(a_2,\dots,a_k)$ gives $[k-1]_t!\,S^\varphi_{k-1}(\hx_i)$.
\end{proof}

\begin{lemma}[Residue form of the HL kernel]\label{lem:res}
Let $\mathbb{K}\supseteq\mathbb{Q}(t)$ be a field and $H\in\mathbb{K}(v)$ with $H(0)=0$, such that $H(v)/v$ has only simple poles, at points $p\in\mathbb{K}$, and $H(v)/v=O(1/v)$ at $\infty$. Then in $\mathbb{K}(x_1,\dots,x_m)$
\[ (1-t)\sum_{i=1}^mH(x_i)\prod_{j\neq i}a_{ij}=-\sum_p\Res_{v=p}f-\Res_{v=\infty}f,\qquad f(v):=\frac{H(v)}{v}\prod_{j=1}^m\frac{v-tx_j}{v-x_j}. \]
\end{lemma}
\begin{proof}
$f$ is regular at $v=0$, its residue at $v=x_i$ is $(1-t)H(x_i)\prod_{j\neq i}a_{ij}$, and the sum of all residues of a rational function is zero.
\end{proof}
For $H(v)=v^n$, $n\ge1$, this is $(1-t)\sum_ix_i^n\prod_{j\ne i}a_{ij}=q_n$ \cite[III (2.10)]{Macdonald95}. Note that $\prod_j(v-tx_j)/(v-x_j)=Q(1/v)$ and $Q(-a)=E(ta)/E(a)$.

%% ===================================================================
\section{The Pieri package}\label{sec:pieri}
%% ===================================================================

\subsection{The coefficients}
For $j\ge0$ put
\[ \alpha_j:=\prod_{i=1}^{j}\frac{s-t^i}{1-t^i}\quad(\alpha_{-1}:=0),\qquad c(n,j):=\frac{(s;t)_{n-j}}{(t;t)_{n-j}}\bigl(\alpha_j-s\,t^{\,n-j}\alpha_{j-1}\bigr)\quad(0\le j\le n), \]
and $c(n,j):=0$ otherwise. Let $F_n(w):=\sum_{j=0}^nc(n,j)w^j$. For instance $F_0=1$ and $F_1(w)=(1-s)(1-tw)/(1-t)$.

\begin{lemma}\label{lem:Cj}
Let $G(x):=\sum_{n\ge0}\frac{(s;t)_n}{(t;t)_n}x^n$, $C_j(x):=\sum_nc(n,j)x^n$, and $D_j:=\alpha_j-s\alpha_{j-1}$. Then
\[ (1-x)G(x)=(1-sx)G(tx),\qquad C_j(x)=D_j\,x^j\,G(tx)\,\frac{1-st^{-j}x}{1-x}, \]
$D_0=1$, $D_j=\alpha_{j-1}t^j(s-1)/(1-t^j)\neq0$ for $j\ge1$, and $(1-t^j)D_j=(ts-t^j)D_{j-1}$ for $j\ge1$.
\end{lemma}
\begin{proof}
Write $g_n=(s;t)_n/(t;t)_n$. The first identity says $g_n-g_{n-1}=t^ng_n-st^{n-1}g_{n-1}$, i.e.\ $(1-t^n)g_n=(1-st^{n-1})g_{n-1}$. By definition $C_j(x)=x^j\bigl(\alpha_jG(x)-s\alpha_{j-1}G(tx)\bigr)=x^jG(tx)\bigl(\alpha_j(1-sx)-s\alpha_{j-1}(1-x)\bigr)/(1-x)$. The bracket is linear in $x$. It equals $D_j$ at $x=0$, and at $x=t^j/s$ it equals $\alpha_j(1-t^j)-\alpha_{j-1}(s-t^j)=0$ (for $j=0$ use $\alpha_{-1}=0$). So it equals $D_j(1-st^{-j}x)$. The formula for $D_j$ follows from $\alpha_j=\alpha_{j-1}(s-t^j)/(1-t^j)$, and the recursion from it and the same relation one step down (for $j=1$ use $D_0=1$).
\end{proof}

\subsection{Statements}
Fix $\ell\ge1$ and indeterminates $z_1,\dots,z_\ell$. For integers $i,j\ge0$ let
\[ K_{ij}(z,w):=\prod_{p=0}^{i-1}\frac{t^pz-sw}{t^pz-t^jw}\ \prod_{r=0}^{j-1}\frac{sz-t^rw}{t^iz-t^rw}, \]
and note $K_{ij}(z,w)=K_{ji}(w,z)$. For $I=(i_1,\dots,i_\ell)\in\mathbb{Z}_{\ge0}^\ell$ with $|I|=S$ put $K_I:=\prod_{c<c'}K_{i_ci_{c'}}(z_c,z_{c'})$ and
\[ V^{(n)}_I:=K_I\,s^{(\ell-1)(n-S)}\sum_{n_1+\dots+n_\ell=n}\ t^{-\sum_{c\neq c'}(n_c-i_c)i_{c'}}\ \prod_{c=1}^\ell t^{-n_ci_c}\,c(n_c,i_c)\,z_c^{-n_c}, \]
with $V^{(n)}_I:=0$ if some $i_c<0$. Since $c(n_c,i_c)=0$ for $n_c<i_c$, $V^{(n)}_I=0$ unless $S\le n$.

\begin{theorem}[$\ell$-column rule]\label{thm:ellcol}
For all $m,k\ge0$ and $\ell\ge1$, in $\Lam_m\otimes\mathbb{Q}(s,t)(z_1,\dots,z_\ell)$,
\[ \mathcal E_k:=\sum_{a\in\mathbb{Z}_{\ge0}^\ell}z_1^{a_1}\cdots z_\ell^{a_\ell}\ e_k\star(e_{a_1}\cdots e_{a_\ell})=\sum_{b=0}^{k}\ \sum_{I}\ \bigl(s^\ell t^{-|I|}\bigr)^b\,V^{(k-b)}_I\ e_b\prod_{c=1}^\ell E(t^{i_c}z_c). \]
In particular the right-hand side is a polynomial in $z_1,\dots,z_\ell$, and $e_k\star(e_{a_1}\cdots e_{a_\ell})$ is its coefficient of $z^a$. The coefficients do not depend on $m$, so the identity holds in $\Lam$.
\end{theorem}

\begin{corollary}[Pieri rule for $e_k\star e_r$]\label{cor:pieri-ekr}
For all $k\ge1$ and $r\ge0$,
\[ e_k\star e_r=\sum_{b=0}^{k}s^b\,F_{k-b}(t^{r-b})\,e_b\,e_{r+k-b}. \]
Equivalently, $\sum_rz^r\,e_k\star e_r=\sum_{j+b\le k}s^bt^{-kj}c(k-b,j)\,e_b\,z^{b-k}E(t^jz)$.
\end{corollary}
\begin{proof}
Take $\ell=1$ in \cref{thm:ellcol}: $K_I=1$, and $V^{(n)}_{(j)}=t^{-nj}c(n,j)z^{-n}$, so the right-hand side is the stated generating function. Extract $[z^r]$ using $[z^r]z^{b-k}E(t^jz)=t^{j(r+k-b)}e_{r+k-b}$.
\end{proof}
For $r<k$, distinct $b$ may give the same monomial $e_be_{r+k-b}$; the identity is between the sums. Since $\star$ is commutative, applying \cref{cor:pieri-ekr} to $e_r\star e_k$ shows that the support of $e_k\star e_r$ is contained in $\{e_be_{r+k-b}:0\le b\le\min(k,r)\}$ (see also \cref{thm:DS}).

\begin{example}[$k=1$]\label{ex:k1}
$c(1,0)=\frac{1-s}{1-t}$ and $c(1,1)=\frac{t(s-1)}{1-t}$, so $F_1(t^r)=(1-s)[r+1]$ and \cref{cor:pieri-ekr} is Hikita's rule $e_1\star e_r=(1-s)[r+1]e_{r+1}+se_1e_r$ \cite[Thm.~3.12]{Hikita25}.
\end{example}

\begin{example}[$k=2$]\label{ex:k2}
One computes $F_2(w)=\dfrac{(1-s)(1-t^2w)\bigl(1-st+(s-t)w\bigr)}{(1-t)(1-t^2)}$, so for all $r\ge0$
\[ e_2\star e_r=\frac{(1-s)(1-t^{r+2})(1-st+(s-t)t^r)}{(1-t)(1-t^2)}\,e_{r+2}+s(1-s)[r]\,e_1e_{r+1}+s^2e_2e_r . \]
For $r=0$ this is $e_2\star1=e_2$.
\end{example}

\begin{corollary}[The case $t=0$]\label{cor:pieri-t0}
For $k\ge1$, $r\ge0$, put $\mu=\min(k,r)$ and $M=\max(k,r)$. The coefficients of $e_k\star e_r$ are regular at $t=0$, and
\[ (e_k\star e_r)\big|_{t=0}=\sum_{b=0}^{\mu-1}(1-s)s^b\,e_b\,e_{r+k-b}+s^\mu\,e_\mu e_M . \]
\end{corollary}
For $s\in[0,1]$ the weights $(1-s)s^b$ ($b<\mu$) and $s^\mu$ are the law of $\min(G,\mu)$ for a geometric variable $G$ with $P(G=b)=(1-s)s^b$, and the monomials $e_be_{r+k-b}$ ($b<\mu$), $e_\mu e_M$ are pairwise distinct. So at $t=0$, $e_k\star e_r$ is a truncated-geometric average of ordinary products.

\begin{corollary}[Two columns; bounded support]\label{cor:twocol}
For $\ell=2$, \cref{thm:ellcol} reads
\[ \sum_{a,b\ge0}z^aw^b\ e_k\star(e_ae_b)=\sum_{b_0,i,j}s^{2b_0}t^{-b_0(i+j)}\,V^{(k-b_0)}_{ij}\,e_{b_0}E(t^iz)E(t^jw), \]
\[ V^{(n)}_{ij}=K_{ij}(z,w)\sum_{n_1+n_2=n}s^{n-i-j}\,t^{-(n_1-i)j-(n_2-j)i-n_1i-n_2j}\,c(n_1,i)\,c(n_2,j)\,z^{-n_1}w^{-n_2}. \]
In general, $e_k\star(e_{a_1}\cdots e_{a_\ell})$ is a linear combination of products of at most $\ell+1$ elementary symmetric functions.
\end{corollary}
\begin{proof}
The first part is the case $\ell=2$. For the second, expand both sides of \cref{thm:ellcol} in the field of iterated Laurent series $\mathbb{Q}(s,t)((z_\ell))\cdots((z_1))$. The left side is a polynomial, so its expansion is itself. On the right, each term is $R(z)\,e_b\prod_cE(t^{i_c}z_c)$ with $R$ rational and free of $x$, and its coefficient of $z^a$ is a finite $\mathbb{Q}(s,t)$-combination of products $e_b\prod_{c=1}^\ell e_{n_c}$.
\end{proof}

\subsection{The recursion}
By \cref{thm:subset}, $e_k\star(e_{a_1}\cdots e_{a_\ell})=\sum_{|A|=k}c_Ax_A\prod_ce_{a_c}(x_{A^c},sx_A)$, and $\sum_{a}z^a\prod_ce_{a_c}(x_{A^c},sx_A)=\prod_c\prod_{j\notin A}(1+x_jz_c)\prod_{i\in A}(1+sx_iz_c)$. Hence
\begin{equation}\label{eq:Gamma}
\mathcal E_k(x)=\prod_cE(z_c)\cdot S^\varphi_k(x),\qquad \varphi(u):=u\prod_{c=1}^\ell\frac{1+suz_c}{1+uz_c}.
\end{equation}
Since $\varphi(x_i)\prod_cE(z_c)=x_i\prod_c(1+sx_iz_c)\cdot\prod_cE(\hx_i;z_c)$, \cref{lem:peel} gives, for $m\ge1$ and $k\ge1$,
\begin{equation}\label{eq:Rl}
[k]\,\mathcal E_k(x)=\sum_{i=1}^mx_i\prod_{c=1}^\ell(1+sx_iz_c)\prod_{j\neq i}a_{ij}\ \mathcal E_{k-1}(\hx_i),
\end{equation}
where $\mathcal E_{k-1}(\hx_i)$ denotes $\mathcal E_{k-1}$ computed in the $m-1$ variables $\hx_i$. For $m=0$ we have $\mathcal E_0=1$ and $\mathcal E_k=0$ for $k\ge1$.

\subsection{The step map}
Fix $I$ and put $\gamma_c:=t^{i_c}z_c$, $C:=s^\ell t^{-|I|}=s^\ell\prod_cz_c/\prod_c\gamma_c$, and
\[ \mathcal{F}(v):=\frac{\prod_{c=1}^\ell(v-sz_c)}{v\prod_{c=1}^\ell(v-\gamma_c)}=\frac{C}{v}+\sum_{c=1}^\ell\frac{\eta_c}{v-\gamma_c},\qquad \eta_c:=\frac{\prod_{c'}(\gamma_c-sz_{c'})}{\gamma_c\prod_{c'\neq c}(\gamma_c-\gamma_{c'})}. \]

\begin{lemma}[Step map]\label{lem:step}
For $b\ge0$, in $m\ge1$ variables,
\begin{multline*}
(1-t)\sum_{i}x_i\prod_c(1+sx_iz_c)\prod_{j\neq i}a_{ij}\ e_b(\hx_i)\prod_cE(\hx_i;\gamma_c)
=C(1-t^{b+1})\,e_{b+1}\prod_cE(\gamma_c)\\
-\sum_{n=0}^{b}e_n\sum_{c=1}^\ell\eta_c\,\gamma_c^{\,n-b-1}\bigl(E(t\gamma_c)-t^nE(\gamma_c)\bigr)\prod_{c'\neq c}E(\gamma_{c'}).
\end{multline*}
\end{lemma}
\begin{proof}
Write $u=x_i$. Then $e_b(\hx_i)=[y^b]E(y)/(1+uy)$ and $E(\hx_i;\gamma)=E(\gamma)/(1+\gamma u)$, so the left side is $\prod_cE(\gamma_c)\cdot[y^b]\,E(y)\,(1-t)\sum_iH(x_i;y)\prod_{j\ne i}a_{ij}$ with
\[ H(v;y):=\frac{v\prod_c(1+svz_c)}{(1+yv)\prod_c(1+\gamma_cv)} \]
(the sum over $i$ is finite and each term is regular at $y=0$, so $[y^b]$ commutes with it). Apply \cref{lem:res} over $\mathbb{K}=\mathbb{Q}(s,t,z,y)$. The poles of $H/v$ are simple, at $v=-1/a$ for $a\in\{y,\gamma_1,\dots,\gamma_\ell\}$, and $H(v)/v\sim C/(yv)$ at $\infty$, which contributes $C/y$. At $v=-1/a$ the residue of $f$ is $\rho_a\,E(ta)/E(a)$, where
\[ \rho_a=\frac{\prod_c(a-sz_c)}{a\prod_{a'\neq a}(a-a')} , \]
as one sees from $\prod_c(1-sz_c/a)=a^{-\ell}\prod_c(a-sz_c)$, $\prod_{a'\neq a}(1-a'/a)=a^{-\ell}\prod_{a'\ne a}(a-a')$ and $Q(-a)=E(ta)/E(a)$. Thus $\rho_y=\mathcal{F}(y)$ and $\rho_{\gamma_c}=-\eta_c/(y-\gamma_c)$, and
\[ (1-t)\sum_iH(x_i;y)\prod_{j\ne i}a_{ij}=\frac Cy-\mathcal F(y)\frac{E(ty)}{E(y)}+\sum_c\frac{\eta_c}{y-\gamma_c}\,\frac{E(t\gamma_c)}{E(\gamma_c)} . \]
Multiply by $E(y)\prod_cE(\gamma_c)$ and write $E(y)=\sum_ne_ny^n$, $E(ty)=\sum_nt^ne_ny^n$. The coefficient of $e_n$ is $[y^{b-n}]$ of the Laurent expansion at $y=0$ of
\begin{multline*} \Bigl(\frac Cy-t^n\mathcal F(y)\Bigr)\prod_cE(\gamma_c)+\sum_c\frac{\eta_c}{y-\gamma_c}E(t\gamma_c)\prod_{c'\ne c}E(\gamma_{c'})\\=\frac{C(1-t^n)}{y}\prod_cE(\gamma_c)+\sum_c\frac{\eta_c}{y-\gamma_c}\bigl(E(t\gamma_c)-t^nE(\gamma_c)\bigr)\prod_{c'\ne c}E(\gamma_{c'}), \end{multline*}
by the partial fraction expansion of $\mathcal F$. Expanding $1/(y-\gamma)=-\sum_{p\ge0}y^p\gamma^{-p-1}$ gives the claim: $n=b+1$ picks the pole at $0$, $n\le b$ picks the second sum, and $n>b+1$ contributes nothing.
\end{proof}

\subsection{Proof of \cref{thm:ellcol}}
Write $\mathcal T_k$ for the right-hand side of \cref{thm:ellcol}, so $\mathcal T_k=\sum_{b,I}W^{(k)}_{b,I}\,e_b\prod_cE(t^{i_c}z_c)$ with $W^{(k)}_{b,I}=C_I^{\,b}V^{(k-b)}_I$, $C_I:=s^\ell t^{-|I|}$; $\epsilon_1,\dots,\epsilon_\ell$ are the unit vectors of $\mathbb{Z}^\ell$.

\emph{Reduction to an index identity.} We show the step identity
\begin{equation}\label{eq:Ll}
\sum_ix_i\prod_c(1+sx_iz_c)\prod_{j\ne i}a_{ij}\,\mathcal T_{k-1}(\hx_i)=[k]\,\mathcal T_k(x)\qquad(m\ge0,\ k\ge1)
\end{equation}
(for $m=0$ the left side is empty). Given \eqref{eq:Ll}, the theorem follows by induction on $k$ for all $m$ at once: $\mathcal E_0=\prod_cE(z_c)=\mathcal T_0$ because $V^{(0)}_{0}=1$; for $m=0$, \eqref{eq:Ll} says $\mathcal T_k=0=\mathcal E_k$; for $m\ge1$, \eqref{eq:Rl}, the induction hypothesis in the variables $\hx_i$, and \eqref{eq:Ll} give $[k]\mathcal E_k=[k]\mathcal T_k$.

Apply \cref{lem:step} to each term of $\mathcal T_{k-1}(\hx_i)$. The term $e_{b'}\prod_cE(t^{i_c}z_c)$ on the right receives three kinds of contributions: from $(b'-1,I)$ via the pole at $0$; from $(b,I)$, $b\ge b'$, via $-t^{b'}E(\gamma_c)$; and from $(b,I-\epsilon_c)$, $b\ge b'$, via $E(t\gamma'_c)$ with $\gamma'_c=t^{i_c-1}z_c$. Writing $\eta'_c$ for $\eta_c$ computed at $I-\epsilon_c$, $n:=k-b'$, $p:=b-b'$, and using $C_{I-\epsilon_c}=tC_I$, the coefficient of $e_{b'}\prod_cE(t^{i_c}z_c)$ in $(1-t)$ times the left side of \eqref{eq:Ll} is $C_I^{\,b'}$ times
\[ (1-t^{b'})V^{(n)}_I+t^{b'}\sum_{p\ge0}C_I^{\,p}\Bigl[V^{(n-1-p)}_I\sum_c\eta_c\gamma_c^{-p-1}-t^p\sum_cV^{(n-1-p)}_{I-\epsilon_c}\eta'_c\,(t^{i_c-1}z_c)^{-p-1}\Bigr]. \]
So \eqref{eq:Ll} follows from the index identity
\begin{equation}\label{eq:starl}
\sum_{p\ge0}C_I^{\,p}\Bigl[V^{(n-1-p)}_I\sum_c\eta_c\gamma_c^{-p-1}-t^p\sum_cV^{(n-1-p)}_{I-\epsilon_c}\eta'_c\,(t^{i_c-1}z_c)^{-p-1}\Bigr]=(1-t^n)V^{(n)}_I\qquad(n\ge0),
\end{equation}
because then the coefficient is $C_I^{\,b'}(1-t^{b'+n})V^{(n)}_I=(1-t^k)W^{(k)}_{b',I}$. (No uniqueness of such expansions is needed: the two sides agree term by term.)

\emph{Generating function.} Let $\mathcal V_I(x):=\sum_{n\ge0}V^{(n)}_Ix^n$. The $p$-sums in \eqref{eq:starl} are convolutions, and $\sum_{p\ge0}C^p\gamma^{-p-1}x^{p+1}=x/(\gamma-Cx)$. So \eqref{eq:starl} for all $n$ is equivalent to
\begin{equation}\label{eq:starlGF}
\mathcal V_I(x)\,U(x)-\mathcal V_I(tx)=-\sum_c\eta'_c\,\frac{x}{t^{i_c-1}z_c-C_Itx}\,\mathcal V_{I-\epsilon_c}(x),\qquad U(x):=1-\sum_c\frac{\eta_cx}{\gamma_c-C_Ix}.
\end{equation}
Put $X_c:=C_Ix/(sz_c)$ and $A_c:=t^{i_c}$. Summing over $(n_c)$ in the definition of $V^{(n)}_I$, the exponent of $t$ attached to column $c$ is $-n_c|I|+i_c(|I|-i_c)$, so
\begin{equation}\label{eq:VGF}
\mathcal V_I(x)=K_I\,s^{-(\ell-1)|I|}\,t^{|I|^2-\sum_ci_c^2}\prod_cC_{i_c}(X_c),
\end{equation}
and passing from $I$ to $I-\epsilon_c$ replaces every $X_{c'}$ by $tX_{c'}$.

\emph{Step A: $U$ is a product.} By the partial fraction expansion of $\mathcal F$, $U(x)=x\,\mathcal F(C_Ix)$. Since $C_Ix-sz_c=-sz_c(1-X_c)$, $C_Ix-\gamma_c=-\gamma_c(1-sX_c/A_c)$ and $C_I\prod_c\gamma_c=s^\ell\prod_cz_c$,
\[ U(x)=\prod_{c=1}^\ell\frac{1-X_c}{1-sX_c/A_c}. \]

\emph{Step B: normalisation.} Divide \eqref{eq:starlGF} by $\mathcal N:=K_I\,s^{-(\ell-1)|I|}t^{|I|^2-\sum i_c^2}\prod_cD_{i_c}X_c^{i_c}G(t^2X_c)$, which is legitimate since $D_j\ne0$. By \cref{lem:Cj} and $G(tX)/G(t^2X)=(1-stX)/(1-tX)$, with $u_c:=stX_c$:
\begin{itemize}
\item column $c$ of $\mathcal V_IU$ contributes $\frac{(1-stX_c)(1-sX_c/A_c)}{(1-tX_c)(1-X_c)}\cdot\frac{1-X_c}{1-sX_c/A_c}=\frac{1-u_c}{1-tX_c}$, and column $c$ of $\mathcal V_I(tx)$ contributes $\frac{A_c-u_c}{1-tX_c}$;
\item in the $c$-th term on the right (present only if $i_c\ge1$), the columns $c'\ne c$ contribute $\frac{A_{c'}-u_{c'}}{1-tX_{c'}}$; column $c$ contributes $\frac{D_{i_c-1}}{D_{i_c}}\frac{A_c}{t}X_c^{-1}\frac{1-st^2X_c/A_c}{1-tX_c}$ with $D_{i-1}/D_i=(1-A)/(ts-A)$; the prefactors contribute $s^{\ell-1}t^{-2(|I|-i_c)}$; and $\frac{x}{t^{i_c-1}z_c-C_Itx}=\frac{stX_c}{C_I(A_c-st^2X_c)}$.
\end{itemize}
It remains to compute $\Phi_c:=\eta'_cK_{I-\epsilon_c}/K_I$. The diagonal factor of $\eta'_c$ is $(\gamma'_c-sz_c)/\gamma'_c=(A_c-st)/A_c$. For $c'\neq c$ put $z=z_c$, $w=z_{c'}$, $i=i_c$, $j=i_{c'}$, using $K_{ij}(z,w)=K_{ji}(w,z)$ to place $c$ in the first slot. Then
\[ \frac{t^{i-1}z-sw}{t^{i-1}z-t^jw}\cdot\frac{K_{i-1,j}(z,w)}{K_{ij}(z,w)}=\prod_{r=0}^{j-1}\frac{t^iz-t^rw}{t^{i-1}z-t^rw}=t^j\,\frac{t^iz-w}{t^iz-t^jw}, \]
the last product telescoping after writing $t^iz-t^rw=t(t^{i-1}z-t^{r-1}w)$. With $\rho_{cc'}:=z_c/z_{c'}=X_{c'}/X_c$ we get
\[ \Phi_c=\frac{A_c-st}{A_c}\prod_{c'\neq c}A_{c'}\frac{A_c\rho_{cc'}-1}{A_c\rho_{cc'}-A_{c'}} . \]
Multiplying everything, $\frac{stX_c}{C_I(A_c-st^2X_c)}\cdot\frac{A_c}{t}X_c^{-1}(1-st^2X_c/A_c)=s/C_I$, then $\frac{s}{C_I}s^{\ell-1}t^{-2(|I|-i_c)}=t^{2i_c-|I|}=A_c/\prod_{c'\ne c}A_{c'}$, and $(A_c-st)/(ts-A_c)=-1$. Hence, after division by $\mathcal N$ and multiplication by $\prod_c(1-tX_c)$, \eqref{eq:starlGF} becomes
\begin{equation}\label{eq:closing}
\prod_c(1-u_c)-\prod_c(A_c-u_c)-\sum_c(1-A_c)\prod_{c'\ne c}\frac{A_cX_{c'}-X_c}{A_cX_{c'}-A_{c'}X_c}\,(A_{c'}-u_{c'})=0 .
\end{equation}
(For $i_c=0$ the $c$-th term is absent, consistently with the factor $1-A_c=0$.)

\emph{Step C: the closing identity.} Put $\lambda_c:=A_c/u_c$ and $\mu_c:=1/u_c$. Since $X_c\propto u_c$, $\frac{A_cX_{c'}-X_c}{A_cX_{c'}-A_{c'}X_c}=\frac{\lambda_c-\mu_{c'}}{\lambda_c-\lambda_{c'}}$, and $(A_{c'}-u_{c'})/u_{c'}=\lambda_{c'}-1$, $(1-A_c)/u_c=\mu_c-\lambda_c$. Dividing \eqref{eq:closing} by $\prod_cu_c$ turns it into the case of \cref{lem:Z} below. As rational functions of $x$, the $\lambda_c=t^{i_c}z_c/(tC_Ix)$ are pairwise distinct and different from $1$. Hence \eqref{eq:starlGF} holds, so \eqref{eq:starl} holds for every $n$, which proves \eqref{eq:Ll} and the theorem. \qed

\begin{lemma}\label{lem:Z}
For indeterminates $\lambda_1,\dots,\lambda_\ell,\mu_1,\dots,\mu_\ell$,
\[ \prod_c(\mu_c-1)-\prod_c(\lambda_c-1)=\sum_{c=1}^\ell(\mu_c-\lambda_c)\prod_{c'\ne c}(\lambda_{c'}-1)\frac{\lambda_c-\mu_{c'}}{\lambda_c-\lambda_{c'}} . \]
\end{lemma}
\begin{proof}
The rational function $g(w)=\prod_c(w-\mu_c)/\bigl((w-1)\prod_c(w-\lambda_c)\bigr)$ has residue $\frac{\lambda_c-\mu_c}{\lambda_c-1}\prod_{c'\ne c}\frac{\lambda_c-\mu_{c'}}{\lambda_c-\lambda_{c'}}$ at $\lambda_c$, residue $\prod_c(1-\mu_c)/\prod_c(1-\lambda_c)$ at $1$, and residue $-1$ at $\infty$. The residues sum to zero; multiply by $\prod_c(\lambda_c-1)$.
\end{proof}

\begin{remark}
For $\ell=1$, \cref{lem:Z} is the tautology $(\mu-1)-(\lambda-1)=\mu-\lambda$, and the proof above reduces to a one-index recursion for $c(n,j)$. For $\ell=2$ it is the polynomial identity that closes the two-column case. The cross kernel $K_I$ is forced by the pairwise splitting of $\eta'_c$ in Step B.
\end{remark}

\subsection{Proof of \cref{cor:pieri-t0}}
At $t=0$, $(t;t)_n=1$, so every $\alpha_j$ and $c(n,j)$ is regular at $t=0$; $(s;t)_n|_{t=0}=1-s$ for $n\ge1$; hence $F_n(0)=c(n,0)|_{t=0}=1-s$ for $n\ge1$, and $F_0=1$.

\emph{Case $r\ge k$.} Every $b\in[0,k]$ has $r-b\ge0$, so each term $s^bF_{k-b}(t^{r-b})=s^b\sum_jc(k-b,j)t^{j(r-b)}$ of \cref{cor:pieri-ekr} is regular at $t=0$. If $b<r$, the terms with $j\ge1$ vanish at $t=0$, leaving $(1-s)s^b$ for $b<k$ and $s^k$ for $b=k$. If $b=r$, then $b=k=r$ and the term is $s^kF_0(1)=s^k$. Summing gives the claim with $\mu=k$.

\emph{Case $r<k$.} By commutativity of $\star$ \cite{Hikita25}, $e_k\star e_r=e_r\star e_k$. For $r=0$ the claim is $e_k\star1=e_k$; for $r\ge1$ apply the first case to $(r,k)$. \qed

\noindent In the $(k,r)$-expansion of \cref{cor:pieri-ekr} with $r<k$, the terms with $b>r$ carry negative powers $t^{j(r-b)}$; the argument shows that these cancel once each monomial $e_be_{r+k-b}$ is collected with its partner $b'=r+k-b$.

\begin{remark}[Other forms and earlier cases]
A second, independent proof of the Pieri rule, in the form $e_k\star e_r=\sum_{y}s^y\mathcal C_{k-y,r-y}e_{k+r-y}e_y$ with $\mathcal C_{a,b}=(-1)^b\frac{[a+b]}{[a]}\pi_a(b)$, is given in \cref{sec:box} (\cref{cor:pieri}); it goes through the plethystic linear coefficient of \cref{thm:lin} and the Column Lemma. The two forms agree after collecting terms. The kernel form of $e_k(Y)$ on symmetric functions is classical in spirit (Macdonald's $D$-operators); what is new here is the level-one transfer $\pi^k$, the explicit coefficients, and the $\ell$-column rule.
\end{remark}

%% ===================================================================
\section{Dominance support}\label{sec:DS}
%% ===================================================================
For a partition $\lambda$ write
\[ e^\star_\lambda:=e_{\lambda_1}\star\cdots\star e_{\lambda_\ell}=E_{\lambda_1}\cdots E_{\lambda_\ell}(1)=\sum_\mu c_{\lambda\mu}(s,t)\,e_\mu , \]
with $E_k$ as in \cref{thm:subset}. The proofs below apply the operators $E_{\lambda_i}$ in an arbitrary order, so they use neither commutativity nor associativity of $\star$.

\begin{theorem}[Dominance support; the $s=0$ valuation]\label{thm:DS}
$e^\star_\lambda=s^{n(\lambda)}e_\lambda+\sum_{\mu\triangleright\lambda}c_{\lambda\mu}e_\mu$ with $c_{\lambda\mu}\in\mathbb{Q}[s,t]$ and $c_{\lambda\mu}|_{s=1}=0$. Moreover $c_{\lambda\mu}\neq0$ for every $\mu\trianglerighteq\lambda$, and the $s$-adic valuation of $c_{\lambda\mu}$ is exactly $n(\mu)$. At $t=0$ the $s$-leading coefficients $[s^{n(\mu)}]c_{\lambda\mu}$ are all $1$: in the basis $b_\mu=s^{n(\mu)}e_\mu$, $e^\star_\lambda$ at $s=t=0$ is the zeta function of dominance order.
\end{theorem}

\begin{corollary}[Operator form]\label{cor:opDS}
For $k\ge1$ and every partition $\mu$, $e_k\star e_\mu\in s^{\sum_i\min(\mu_i,k)}e_{\mu\cup k}+\mathrm{span}\{e_\nu:\nu\triangleright\mu\cup k\}$.
\end{corollary}

We work in $\Lam_m$ with $m\ge n=|\lambda|$; by \cref{lem:stable} this loses nothing.

\subsection{Degrees and dominance}
\begin{lemma}\label{lem:e-m}
Let $m\ge n$ and $\rho\vdash n$. Then $\mathrm{span}\{e_\nu:\nu\trianglerighteq\rho'\}=\mathrm{span}\{m_\kappa:\kappa\trianglelefteq\rho\}$ in the degree-$n$ part of $\Lam_m$.
\end{lemma}
\begin{proof}
By \cite[I (6.6), (6.7)]{Macdonald95}, $e_\nu=\sum_\kappa M_{\nu\kappa}m_\kappa$, where $M_{\nu\kappa}$ is the number of $0$-$1$ matrices with row sums $\nu$ and column sums $\kappa$; $M_{\nu\kappa}\ne0$ forces $\kappa\trianglelefteq\nu'$, and $M_{\nu\nu'}=1$. Since $\nu\trianglerighteq\rho'\iff\nu'\trianglelefteq\rho$ \cite[I (1.11)]{Macdonald95}, the left side is contained in the right. Both have dimension $\#\{\kappa\vdash n:\kappa\trianglelefteq\rho\}$, because the $e_\nu$, $\nu\vdash n$, are independent for $m\ge n$.
\end{proof}

For $S\subseteq[m]$ and nonzero $R\in\mathbb{Q}(s,t)(x)$, let $\deg_SR$ be the degree in $u$ of $R|_{x_i\mapsto ux_i\,(i\in S)}$ (numerator degree minus denominator degree). It is additive on products and satisfies $\deg_S(R_1+R_2)\le\max(\deg_SR_1,\deg_SR_2)$. For a polynomial $P$, $\deg_SP$ is the largest value of $\sum_{i\in S}\alpha_i$ over the monomials $x^\alpha$ of $P$; for symmetric $P$ it depends only on $j=|S|$, and we write $D_j(P)$. Since the $m_\kappa$ have disjoint supports and $D_j(m_\kappa)=\kappa_1+\dots+\kappa_j$:

\begin{lemma}\label{lem:Dj}
A homogeneous $P\in\Lam_m$ of degree $|\rho|$ lies in $\mathrm{span}\{m_\kappa:\kappa\trianglelefteq\rho\}$ if and only if $D_j(P)\le\rho_1+\dots+\rho_j$ for all $j$.
\end{lemma}

\begin{lemma}[Degree step]\label{lem:degstep}
For symmetric $F$ and every $j$, $D_j(E_kF)\le D_j(F)+\min(k,j)$.
\end{lemma}
\begin{proof}
Fix $|S|=j$ and a term $R_A=c_A\,x_A\,F|_{x_i\mapsto sx_i\,(i\in A)}$ of \cref{thm:subset}. Every $a_{ij}$ has $\deg_S a_{ij}=0$: if $i,j$ are both in $S$ or both outside, $a_{ij}$ is homogeneous of degree $0$ in $u$ or constant; otherwise numerator and denominator have degree $1$. Next, $\deg_Sx_A=|A\cap S|\le\min(k,j)$. Rescaling some variables by the constant $s$ does not change $u$-degrees, so the last factor has $\deg_S=D_j(F)$. Since $E_kF$ is a polynomial, its $\deg_S$ is its monomial degree.
\end{proof}

\subsection{Support, diagonal coefficient, and the value at $s=1$}
\begin{proof}[Proof of the support statement]
Starting from $D_j(1)=0$ and applying \cref{lem:degstep} once for each part, $D_j(e^\star_\lambda)\le\sum_i\min(\lambda_i,j)=\lambda'_1+\dots+\lambda'_j$. By \cref{lem:Dj,lem:e-m}, $e^\star_\lambda\in\mathrm{span}\{m_\kappa:\kappa\trianglelefteq\lambda'\}=\mathrm{span}\{e_\nu:\nu\trianglerighteq\lambda\}$. The support part of \cref{cor:opDS} is the same argument with one step, starting from $D_j(e_\mu)=\sum_i\min(\mu_i,j)$.
\end{proof}

Fix integers $w_1>w_2>\dots>w_m>0$ and for $\rho=(\rho_1\ge\dots\ge\rho_m\ge0)$ put $W(\rho)=\sum_iw_i\rho_i$.

\begin{lemma}\label{lem:weights}
If $\mathrm{sort}(\alpha)=\kappa\trianglelefteq\rho$, then $\sum_iw_i\alpha_i\le W(\rho)$, with equality only if $\alpha=\rho$. Hence if $P\in\mathrm{span}\{m_\kappa:\kappa\trianglelefteq\rho\}$ has $m_\rho$-coefficient $c$, the coefficient of $u^{W(\rho)}$ in $P(u^{w_1}x_1,\dots,u^{w_m}x_m)$ is $c\,x^\rho$, and no higher power of $u$ occurs.
\end{lemma}
\begin{proof}
$\sum_iw_i\alpha_i\le\sum_iw_i\kappa_i$ by the rearrangement inequality, with equality only if $\alpha=\kappa$ (the $w_i$ are distinct). By Abel summation, $\sum_iw_i\kappa_i=\sum_j(w_j-w_{j+1})(\kappa_1+\dots+\kappa_j)$ with $w_{m+1}:=0$; all weights are positive, so this is at most the same expression for $\rho$, with equality only if $\kappa=\rho$.
\end{proof}

\begin{lemma}[Lead step]\label{lem:leadstep}
Let $F\in\mathrm{span}\{m_\kappa:\kappa\trianglelefteq\rho\}$ with $m_\rho$-coefficient $c_F$. Then $E_kF\in\mathrm{span}\{m_\kappa:\kappa\trianglelefteq\rho+1^k\}$, with $m_{\rho+1^k}$-coefficient $s^{\rho_1+\dots+\rho_k}c_F$.
\end{lemma}
\begin{proof}
Membership is \cref{lem:degstep,lem:Dj}. Substitute $x_i\mapsto u^{w_i}x_i$ and expand each term $R_A$ in $u^{-1}$. For $i<j$, $a_{ij}=1+O(u^{-1})$, and for $i>j$, $a_{ij}=t+O(u^{-1})$, so $c_A=t^{\mathrm{inv}(A)}(1+O(u^{-1}))$ with $\mathrm{inv}(A)=\#\{i\in A,j\notin A:i>j\}$. Next, $x_A$ becomes $u^{\sum_{a\in A}w_a}x_A$. Finally, $F|_{x_a\mapsto sx_a\,(a\in A)}$ has the monomials of $F$ with $x^\alpha$ rescaled by $s^{\sum_{a\in A}\alpha_a}$, so by \cref{lem:weights} it is $u^{W(\rho)}\bigl(c_Fs^{\sum_{a\in A}\rho_a}x^\rho+O(u^{-1})\bigr)$. The exponent $W(\rho)+\sum_{a\in A}w_a$ is uniquely maximal for $A=[k]$, where $\mathrm{inv}=0$. So the coefficient of $u^{W(\rho+1^k)}$ in $E_kF$ is $s^{\rho_1+\dots+\rho_k}c_F\,x^{\rho+1^k}$; conclude by \cref{lem:weights} for $\rho+1^k$.
\end{proof}

\begin{proof}[Proof of the diagonal coefficient]
Apply the operators in any order. If the parts applied so far form $\mu$, the current function lies in $\mathrm{span}\{m_\kappa:\kappa\trianglelefteq\mu'\}$, and applying $E_k$ multiplies the $m_{\mu'}$-coefficient by $s^{\mu'_1+\dots+\mu'_k}=s^{\sum_i\min(\mu_i,k)}$ (\cref{lem:leadstep}; note $\mu'+1^k=(\mu\cup k)'$). Each unordered pair of parts contributes its minimum once, so the total exponent is $\sum_{i<p}\min(\lambda_i,\lambda_p)=\sum_p(p-1)\lambda_p=n(\lambda)$. By the triangularity in the proof of \cref{lem:e-m}, the only $e_\nu$ with $\nu\trianglerighteq\lambda$ containing $m_{\lambda'}$ is $e_\lambda$, with coefficient $1$. Hence $c_{\lambda\lambda}=s^{n(\lambda)}$. The same computation gives the leading term in \cref{cor:opDS}.
\end{proof}

\begin{proof}[Proof of integrality and of $c_{\lambda\mu}|_{s=1}=0$]
Let $V=\prod_{i<j}(x_i-x_j)$. Multiplying the subset formula by $V$ gives
\[ V\,E_kF=\sum_{|A|=k}\varepsilon_A\prod_{\substack{i<j\\ \text{both in or both out of }A}}(x_i-x_j)\prod_{i\in A,\,j\notin A}(x_i-tx_j)\ x_A\,F|_{x_a\mapsto sx_a\,(a\in A)} ,\qquad\varepsilon_A=\pm1. \]
If $F\in\mathbb{Q}[s,t][x]$, the right side lies in $\mathbb{Q}[s,t][x]$. As $V$ has content $1$ and $E_kF$ is a polynomial, Gauss's lemma gives $E_kF\in\mathbb{Q}[s,t][x]$. By induction $e^\star_\lambda\in\mathbb{Q}[s,t][x]^{S_m}$, and its $e$-coefficients lie in $\mathbb{Q}[s,t]$. Now set $s=1$ (a ring map on $\mathbb{Q}[s,t][x]$): every $F|_{x_a\mapsto sx_a}$ becomes $F|_{s=1}$, so $V\,(E_kF)|_{s=1}=F|_{s=1}\cdot V\,E_k(1)=V\,e_k\,F|_{s=1}$ by \cref{lem:Ek1}. Hence $(E_kF)|_{s=1}=e_kF|_{s=1}$, $e^\star_\lambda|_{s=1}=e_\lambda$, and $c_{\lambda\mu}|_{s=1}=\delta_{\lambda\mu}$.
\end{proof}

\subsection{The exact $s$-adic valuation}
It remains to show that $s^{n(\mu)}$ divides $c_{\lambda\mu}$ and that $(c_{\lambda\mu}/s^{n(\mu)})|_{s=t=0}=1$ for all $\mu\trianglerighteq\lambda$; in particular $c_{\lambda\mu}\ne0$. Put $\omega(\alpha):=\sum_j\binom{\alpha_j}{2}$ for $\alpha\in\mathbb{Z}^m_{\ge0}$ and $B_\mu:=\{\beta:\mathrm{sort}(\beta)\trianglelefteq\mu'\}$, the monomial support of $e^\star_\mu$. Write $e^\star_\lambda=\sum_\beta a^\lambda_\beta x^\beta$; $a^\lambda_\beta$ is symmetric in $\beta$. The source of $s^{n}$ is the identity
\begin{equation}\label{eq:omega}
\omega(\beta+1_A)=\omega(\beta)+1_A\cdot\beta .
\end{equation}

Let $\mathbb{K}=\mathbb{Q}(t)(s^{1/2})$ (or $\mathbb{Q}(s^{1/2})$ when $t=0$) with the $s$-adic valuation $v$, extended to $\mathbb{K}(x)$ by the Gauss valuation $v(\sum a_\beta x^\beta)=\min_\beta v(a_\beta)$. For $c\in(\frac12\mathbb{Z})^m$ let $\sigma_c$ be the automorphism $x_j\mapsto s^{-c_j}x_j$, so that $v(\sigma_cP)=\min_\beta(v(a_\beta)-c\cdot\beta)$.

\begin{lemma}\label{lem:valA}
For generic $t$, and also at $t=0$, $v(a^\lambda_\alpha)\ge\omega(\alpha)$ for all $\alpha$.
\end{lemma}
\begin{proof}
We show $v(\sigma_ce^\star_\lambda)\ge\min_{\beta\in B_\lambda}(\omega(\beta)-c\cdot\beta)$ for all $c$, by induction along $e^\star_\mu\mapsto E_ke^\star_\mu=e^\star_{\mu\cup k}$, the case $e^\star_\varnothing=1$ being clear. In a term $R_A$: $v(\sigma_ca_{ij})=0$ for generic $t$, and at $t=0$, $v(\sigma_ca_{ij})=\max(c_i,c_j)-c_i\ge0$; $\sigma_cx_A=s^{-c\cdot1_A}x_A$; and $\sigma_c(F|_{x_a\mapsto sx_a\,(a\in A)})=\sigma_{c-1_A}F$. Hence, by induction and \eqref{eq:omega},
\[ v(\sigma_cR_A)\ge-c\cdot1_A+\min_{\beta\in B_\mu}\bigl(\omega(\beta)-(c-1_A)\cdot\beta\bigr)=\min_{\beta\in B_\mu}\bigl(\omega(\beta+1_A)-c\cdot(\beta+1_A)\bigr), \]
and $\beta+1_A\in B_{\mu\cup k}$ because the top-$j$ sums grow by at most $\min(j,k)$ and $\mu'+1^k=(\mu\cup k)'$. Now take $c=\alpha-\frac12\mathbf{1}$. Then $\omega(\beta)-c\cdot\beta=\sum_j\bigl((\beta_j-\alpha_j)^2-\alpha_j^2\bigr)/2$ is uniquely minimised over $\mathbb{Z}^m$ at $\beta=\alpha$, so $v(a^\lambda_\alpha)-c\cdot\alpha\ge v(\sigma_ce^\star_\lambda)\ge\omega(\alpha)-c\cdot\alpha$.
\end{proof}

\begin{proof}[Proof that $s^{n(\mu)}\mid c_{\lambda\mu}$]
Write $c_{\lambda\nu}=s^{n(\nu)}d_\nu$ and suppose $\delta:=\min_\nu v(d_\nu)<0$. Choose $\mu$ dominance-minimal with $v(d_\mu)=\delta$. The coefficient of $x^{\mu'}$ is $a^\lambda_{\mu'}=\sum_{\nu\trianglelefteq\mu}c_{\lambda\nu}M_{\nu\mu'}$ with $M_{\mu\mu'}=1$. Since $\nu\triangleleft\mu$ implies $n(\nu)>n(\mu)$ (as $n(\nu)=\sum_{i\ge1}(n-\nu_1-\dots-\nu_i)$), every term with $\nu\ne\mu$ has valuation $>n(\mu)+\delta$, so $v(a^\lambda_{\mu'})=n(\mu)+\delta<n(\mu)=\omega(\mu')$, contradicting \cref{lem:valA}. So $d_\nu$ is regular at $s=0$; as $c_{\lambda\nu}\in\mathbb{Q}[s,t]$, $s^{n(\nu)}$ divides $c_{\lambda\nu}$ in $\mathbb{Q}[s,t]$. The same argument at $t=0$ shows $d_\mu|_{s=t=0}=[s^{\omega(\mu')}]\,a^\lambda_{\mu'}|_{t=0}$.
\end{proof}

For the value at $s=t=0$ we need the initial forms. For $R\in\mathbb{K}(x)$ with $v(R)\ge v_0$, let $\mathrm{in}_{v_0}(R)$ be the image of $s^{-v_0}R$ in the residue field $\mathbb{Q}(x)$; it is additive at a common level and multiplicative. For $\kappa=\mathrm{sort}(\alpha)$ let $\mathrm{peel}_k$ subtract $1$ from the $k$ largest parts.

\begin{lemma}[Peel Lemma]\label{lem:peelmono}
If $\kappa\trianglelefteq\rho$, $|\kappa|=|\rho|$ and $k\le\ell(\rho)$, then $\mathrm{peel}_k\kappa\trianglelefteq\mathrm{peel}_k\rho$.
\end{lemma}
\begin{proof}
Note $\ell(\kappa)\ge\ell(\rho)\ge k$. Counting parts $\ge c$ gives $(\mathrm{peel}_k\kappa)'_c=\kappa'_c-\min(k,\kappa'_c)+\min(k,\kappa'_{c+1})$, so with $S_C$ the sum of the first $C$ parts and $\kappa'_1\ge k$, $S_C((\mathrm{peel}_k\kappa)')=S_C(\kappa')-k+\min(k,\kappa'_{C+1})$, and likewise for $\rho$. Let $D_C=S_C(\kappa')-S_C(\rho')\ge0$ (by \cite[I (1.11)]{Macdonald95}) and $\bar m(y)=\min(k,y)$. If $\bar m(\kappa'_{C+1})\ge\bar m(\rho'_{C+1})$ we are done. Otherwise $\kappa'_{C+1}<k$ and $\kappa'_{C+1}<\rho'_{C+1}$, and $D_C=D_{C+1}+\rho'_{C+1}-\kappa'_{C+1}\ge\bar m(\rho'_{C+1})-\bar m(\kappa'_{C+1})$. So $(\mathrm{peel}_k\kappa)'\trianglerighteq(\mathrm{peel}_k\rho)'$.
\end{proof}

\begin{lemma}\label{lem:valB}
At $t=0$, $\Lambda^0_\lambda(\alpha):=[s^{\omega(\alpha)}]\,a^\lambda_\alpha$ equals $1$ if $\mathrm{sort}(\alpha)\trianglelefteq\lambda'$ and $0$ otherwise.
\end{lemma}
\begin{proof}
Induction along $E_k$; the base case is clear, and $\Lambda^0_{\mu\cup k}(\alpha)=0$ outside $B_{\mu\cup k}$ by the support statement. Let $\mathrm{sort}(\alpha)\trianglelefteq(\mu\cup k)'$, $c=\alpha-\frac12\mathbf1$ and $v_0=\omega(\alpha)-c\cdot\alpha$. By the uniqueness in the proof of \cref{lem:valA}, $\mathrm{in}_{v_0}(\sigma_ce^\star_{\mu\cup k})=\Lambda^0_{\mu\cup k}(\alpha)\,x^\alpha$. If $\alpha-1_A$ has a negative entry or lies outside $B_\mu$, then $v(\sigma_cR_A)>v_0$ by that uniqueness and the term $R_A$ does not contribute. Otherwise, with $\beta_0=\alpha-1_A$, the levels split as in \eqref{eq:omega} and
\[ \mathrm{in}_{v_0}(\sigma_cR_A)=\mathrm{in}_0(\sigma_cc_A)\cdot\Lambda^0_\mu(\beta_0)\,x^\alpha . \]
At $t=0$, $\mathrm{in}_0(\sigma_ca_{ij})$ is $1$ if $\alpha_i>\alpha_j$, $0$ if $\alpha_i<\alpha_j$, and $x_i/(x_i-x_j)$ if $\alpha_i=\alpha_j$. So only sets $A$ that are upper sets for $\alpha$ contribute: with $v$ the $k$-th largest entry of $\alpha$ and $L=\{j:\alpha_j=v\}$, $A=\{j:\alpha_j>v\}\sqcup B$ with $B\subseteq L$ of size $r=k-\#\{j:\alpha_j>v\}$. As $\ell(\mathrm{sort}\,\alpha)\ge\ell((\mu\cup k)')\ge k$, $v\ge1$. For all such $A$, $\mathrm{sort}(\alpha-1_A)=\mathrm{peel}_k(\mathrm{sort}\,\alpha)$, and $\Lambda^0_\mu$ is symmetric, so
\[ \Lambda^0_{\mu\cup k}(\alpha)=\Lambda^0_\mu(\mathrm{peel}_k\,\mathrm{sort}\,\alpha)\cdot\sum_{B\subseteq L,\ |B|=r}\ \prod_{i\in B,\ j\in L\setminus B}\frac{x_i}{x_i-x_j}=\Lambda^0_\mu(\mathrm{peel}_k\,\mathrm{sort}\,\alpha). \]
The last sum is the case $t=0$ of $\sum_{|B|=r}\prod_{i\in B,j\in L\setminus B}\frac{x_i-tx_j}{x_i-x_j}=\qbin{|L|}{r}$, which follows from the Poincar\'e identity \cite[III (1.4)]{Macdonald95} by grouping $S_{|L|}$ according to the image of $\{1,\dots,r\}$. Finally, by \cref{lem:peelmono} with $\rho=(\mu\cup k)'=\mu'+1^k$, $\mathrm{peel}_k(\mathrm{sort}\,\alpha)\trianglelefteq\mathrm{peel}_k(\mu'+1^k)=\mu'$, so $\Lambda^0_\mu(\mathrm{peel}_k\,\mathrm{sort}\,\alpha)=1$.
\end{proof}

\begin{proof}[End of proof of \cref{thm:DS}]
By the last two proofs, $(c_{\lambda\mu}/s^{n(\mu)})|_{s=t=0}=\Lambda^0_\lambda(\mu')=1$ for $\mu\trianglerighteq\lambda$, since $\mu'\trianglelefteq\lambda'$. In particular $c_{\lambda\mu}\neq0$ and its $s$-adic valuation is $n(\mu)$.
\end{proof}

\begin{remark}
The leading coefficients $d_{\lambda\mu}(t)=[s^{n(\mu)}]c_{\lambda\mu}$ depend on $t$; for instance $d_{(1^4),(2,1,1)}=1+3t$. They are identified in \cref{cor:dmatrix}. The degree count above is the same kind of triangularity argument as for Macdonald's operators \cite[VI (3.6)--(3.10)]{Macdonald95}.
\end{remark}

%% ===================================================================
\section{The $s\to0$ edge}\label{sec:edge}
%% ===================================================================
Since $s=q^{-1}$, the edge $s\to0$ is Hikita's $q\to\infty$, where the unrescaled product degenerates \cite[Lemma~6.3]{Hikita25}. After the rescaling of \cref{thm:DS} it becomes Hall--Littlewood multiplication. Throughout this section $t$ is an indeterminate, $b_\mu:=s^{n(\mu)}e_\mu$, and we work in $\Lam_m$ with $m$ at least the degree (\cref{lem:stable}). We keep the notation $\mathbb{K}$, $v$, $\sigma_c$, $\omega$ of \cref{sec:DS}, with $\mathbb{K}=\mathbb{Q}(t)(s^{1/2})$.

\begin{theorem}[The $s\to0$ edge]\label{thm:H}
Let $\mathcal O=\mathbb{Q}(t)[s]_{(s)}$ and $\mathcal L=\bigoplus_\mu\mathcal O\,b_\mu$.
\begin{enumerate}
\item $E_kb_\mu\in\bigoplus_\nu\mathbb{Q}[s,t]\,b_\nu$; in particular $E_k\mathcal L\subseteq\mathcal L$.
\item Let $L_k$ be the map induced by $E_k$ on $\mathcal L/s\mathcal L=\bigoplus_\mu\mathbb{Q}(t)\,\bar b_\mu$, and let $\phi(\bar b_\mu):=t^{-n(\mu')}P_{\mu'}(x;t^{-1})$. Then $\phi\circ L_k=t^{-\binom k2}\,e_k\cdot\phi$.
\item Explicitly, with $\rho=\mu'$ and $\kappa=\nu'$,
\[ L_k\bar b_\mu=\sum_{\kappa/\rho\ \text{vertical $k$-strip}}t^{\mathrm{inv}(\kappa/\rho)}\prod_{v\ge1}\qbin{m_v(\kappa)}{r_v}\ \bar b_\nu, \]
where $r_v$ is the number of rows of $\kappa$ of length $v$ that end in a box of the strip, and $\mathrm{inv}(\kappa/\rho)=\sum_{v<w}r_v\,(m_w(\kappa)-r_w)$.
\end{enumerate}
\end{theorem}

\begin{corollary}[The $s$-leading coefficients]\label{cor:dmatrix}
For $\lambda,\mu\vdash n$, $d_{\lambda\mu}(t):=[s^{n(\mu)}]\,c_{\lambda\mu}$ satisfies
\[ d_{\lambda\mu}(t)=t^{\,n(\mu')-n(\lambda')}\,[P_{\mu'}(x;t^{-1})]\,e_\lambda=t^{-n(\lambda')}\sum_\nu K_{\nu'\lambda}\,\widetilde K_{\nu\mu'}(t), \]
where $\widetilde K_{\nu\kappa}(t)=t^{n(\kappa)}K_{\nu\kappa}(t^{-1})$ is the cocharge Kostka--Foulkes polynomial. Hence $d_{\lambda\mu}\in\mathbb{N}[t]$, and $d_{\lambda\mu}(1)$ is the number of $0$-$1$ matrices with row sums $\lambda$ and column sums $\mu'$.
\end{corollary}

\subsection{Admissible functions}
Call a homogeneous $F=\sum_\beta a_\beta x^\beta\in\mathbb{K}[x]^{S_m}$ \emph{admissible} if $v(a_\beta)\ge\omega(\beta)$ for all $\beta$. As in the proof of \cref{lem:valA}, $F$ is admissible if and only if $v(\sigma_cF)\ge\min_{\beta\in\mathbb{Z}^m_{\ge0}}(\omega(\beta)-c\cdot\beta)$ for all $c\in(\frac12\mathbb{Z})^m$; the backward direction uses $c=\alpha-\frac12\mathbf1$, for which $\beta=\alpha$ is the unique minimiser. For admissible $F$ let $\mathrm{in}(F)(\alpha)\in\mathbb{Q}(t)$ be the residue of $s^{-\omega(\alpha)}a_\alpha$; it depends only on $\mathrm{sort}(\alpha)$. With $c=\alpha-\frac12\mathbf1$ and $v_0=\omega(\alpha)-c\cdot\alpha$,
\begin{equation}\label{eq:inform}
\mathrm{in}_{v_0}(\sigma_cF)=\mathrm{in}(F)(\alpha)\,x^\alpha .
\end{equation}

\begin{lemma}\label{lem:admissible}
A homogeneous $F=\sum_\mu c_\mu e_\mu$ of degree $\le m$ is admissible if and only if $v(c_\mu)\ge n(\mu)$ for all $\mu$. In that case $\mathrm{in}(F)(\mu')$ is the residue of $s^{-n(\mu)}c_\mu$.
\end{lemma}
\begin{proof}
The coefficient of $x^\beta$ in $b_\mu$ is $s^{n(\mu)}M_{\mu\kappa}$ with $\kappa=\mathrm{sort}(\beta)$. If $M_{\mu\kappa}\ne0$ then $\kappa\trianglelefteq\mu'$, so $\kappa'\trianglerighteq\mu$ and $\omega(\beta)=n(\kappa')\le n(\mu)$, with equality only for $\kappa=\mu'$, where $M_{\mu\mu'}=1$. So $b_\mu$ is admissible with $\mathrm{in}(b_\mu)(\kappa)=[\kappa=\mu']$, which gives the backward direction and the formula. The forward direction is the argument in the proof that $s^{n(\mu)}\mid c_{\lambda\mu}$ in \cref{sec:DS}, with $a^\lambda$ replaced by the coefficients of $F$.
\end{proof}

\begin{lemma}\label{lem:levelset}
$\sum_{B\subseteq[N],|B|=r}\prod_{i\in B,j\notin B}a_{ij}=\qbin Nr$.
\end{lemma}
\begin{proof}
Group the Poincar\'e identity $\sum_{w\in S_N}w(\prod_{i<j}a_{ij})=[N]_t!$ \cite[III (1.4)]{Macdonald95} by $B=w(\{1,\dots,r\})$; the stabiliser $S_r\times S_{N-r}$ contributes $[r]_t![N-r]_t!$.
\end{proof}

\subsection{The recursion}
\begin{proposition}\label{prop:Hrec}
Let $F$ be admissible of degree $n$, with $m\ge n+k$. Then $E_kF$ is admissible, and for every $\kappa\vdash n+k$
\begin{equation}\label{eq:Hrec}
\mathrm{in}(E_kF)(\kappa)=\sum_{\rho:\ \kappa/\rho\ \text{vertical $k$-strip}}t^{\mathrm{inv}(\kappa/\rho)}\prod_{v\ge1}\qbin{m_v(\kappa)}{r_v}\ \mathrm{in}(F)(\rho).
\end{equation}
\end{proposition}
\begin{proof}
Fix $c$ and a term $R_A$ of \cref{thm:subset}. Since $t$ is a unit, $v(\sigma_ca_{ij})=0$, and the initial form of $\sigma_ca_{ij}$ is $1$ if $c_i>c_j$, $t$ if $c_i<c_j$, and $a_{ij}$ if $c_i=c_j$. As in \cref{lem:valA}, using \eqref{eq:omega},
\[ v(\sigma_cR_A)\ge\min_{\gamma\ge1_A}\bigl(\omega(\gamma)-c\cdot\gamma\bigr)\ge\min_{\gamma\ge0}\bigl(\omega(\gamma)-c\cdot\gamma\bigr), \]
so $E_kF$ is admissible. Now fix $\alpha$ with $\mathrm{sort}(\alpha)=\kappa$, $c=\alpha-\frac12\mathbf1$, $v_0=\omega(\alpha)-c\cdot\alpha$; all terms have $v\ge v_0$, so initial forms add, and by \eqref{eq:inform} $\mathrm{in}(E_kF)(\alpha)x^\alpha=\sum_A\mathrm{in}_{v_0}(\sigma_cR_A)$. If $A\not\subseteq\mathrm{supp}(\alpha)$, then $\gamma=\alpha$ is excluded from $\{\gamma\ge1_A\}$ and $v(\sigma_cR_A)>v_0$. If $A\subseteq\mathrm{supp}(\alpha)$, put $\beta_0=\alpha-1_A$; then $c-1_A=\beta_0-\frac12\mathbf1$, the levels add by \eqref{eq:omega}, and
\[ \mathrm{in}_{v_0}(\sigma_cR_A)=\mathrm{in}_0(\sigma_cc_A)\cdot\mathrm{in}(F)(\beta_0)\,x^\alpha . \]
Let $L_v=\{j:\alpha_j=v\}$ ($v\ge1$), $B_v=A\cap L_v$ and $r_v=|B_v|$. Then
\[ \mathrm{in}_0(\sigma_cc_A)=t^{\#\{i\in A,\,j\notin A:\ \alpha_i<\alpha_j\}}\prod_{v\ge1}\ \prod_{i\in B_v,\,j\in L_v\setminus B_v}a_{ij},\qquad \#\{\cdots\}=\sum_{v<w}r_v(m_w-r_w). \]
The partition $\mathrm{sort}(\beta_0)$ is obtained from $\kappa$ by lowering $r_v$ of its parts equal to $v$, for each $v$; this is the $\rho$ with $\kappa/\rho$ a vertical $k$-strip and data $(r_v)$, and every such strip arises from exactly one $(r_v)$. So both $\mathrm{in}(F)(\beta_0)$ and the power of $t$ depend only on $(r_v)$, and summing over the $B_v$ with \cref{lem:levelset} gives $\prod_v\qbin{m_v}{r_v}$.
\end{proof}
At $t=0$ every $t$-binomial is $1$ and $t^{\mathrm{inv}}$ vanishes unless every part larger than a lowered part is lowered; so \eqref{eq:Hrec} reduces to the recursion in the proof of \cref{lem:valB}.

\subsection{Proof of \cref{thm:H} and \cref{cor:dmatrix}}
\emph{(1).} $b_\mu$ is admissible, so $E_kb_\mu$ is admissible (\cref{prop:Hrec}) and its $b_\nu$-coordinates have $v\ge0$ (\cref{lem:admissible}). By the integrality part of \cref{thm:DS}'s proof, $E_ke_\mu\in\bigoplus_\nu\mathbb{Q}[s,t]e_\nu$, so the $b_\nu$-coordinate is $p/s^{n(\nu)-n(\mu)}$ with $p\in\mathbb{Q}[s,t]$; having $v\ge0$ over $\mathbb{Q}(t)$, it lies in $\mathbb{Q}[s,t]$.

\emph{(3).} By \cref{lem:admissible}, the $\bar b_\nu$-coordinate of $L_k\bar b_\mu$ is $\mathrm{in}(E_kb_\mu)(\nu')$; apply \eqref{eq:Hrec} with $\mathrm{in}(b_\mu)(\rho)=[\rho=\mu']$.

\emph{(2).} Let $\kappa/\rho$ be a vertical $k$-strip with data $(r_v)$, obtained by lowering the parts indexed by $A$. Writing $n(\kappa)=\sum_{\text{pairs of parts}}\min$, a pair of parts both in $A$ loses $1$; a pair with one part of value $v$ in $A$ and one of value $w$ not in $A$ loses $1$ if and only if $v\le w$; other pairs lose nothing. Hence
\begin{equation}\label{eq:ncount}
n(\kappa)-n(\rho)=\binom k2+\mathrm{inv}(\kappa/\rho)+\sum_vr_v(m_v-r_v).
\end{equation}
Since $\qbin mr\big|_{t\mapsto t^{-1}}=t^{-r(m-r)}\qbin mr$, \eqref{eq:ncount} gives $t^{\mathrm{inv}}\prod_v\qbin{m_v}{r_v}=t^{\,n(\kappa)-n(\rho)-\binom k2}\psi_{\kappa/\rho}(t^{-1})$ with $\psi_{\kappa/\rho}(u):=\prod_v\genfrac{[}{]}{0pt}{}{m_v(\kappa)}{r_v}_u$. The Hall--Littlewood Pieri rule \cite[III (3.2)]{Macdonald95} reads $e_kP_\rho(x;u)=\sum_\kappa\psi_{\kappa/\rho}(u)P_\kappa(x;u)$, summed over vertical $k$-strips (there $m_v(\kappa)=\kappa'_v-\kappa'_{v+1}$ and $r_v=\kappa'_v-\rho'_v$). For $X\in\mathcal L$ put $h(\kappa)=t^{-n(\kappa)}\mathrm{in}(X)(\kappa)$, so that $\phi(\bar X)=\sum_\kappa h(\kappa)P_\kappa(x;t^{-1})$ by \cref{lem:admissible}. By \eqref{eq:Hrec}, $t^{-n(\kappa)}\mathrm{in}(E_kX)(\kappa)=t^{-\binom k2}\sum_\rho\psi_{\kappa/\rho}(t^{-1})h(\rho)$, hence $\phi(L_k\bar X)=t^{-\binom k2}\sum_\rho h(\rho)\,e_kP_\rho(x;t^{-1})=t^{-\binom k2}e_k\,\phi(\bar X)$.

\emph{\cref{cor:dmatrix}.} Since $1=b_\varnothing$ and $\phi(\bar b_\varnothing)=1$, iterating (2) gives $\phi(e^\star_\lambda\bmod s)=t^{-\sum_i\binom{\lambda_i}2}e_\lambda=t^{-n(\lambda')}e_\lambda$. By \cref{thm:DS}, $e^\star_\lambda\equiv\sum_\mu d_{\lambda\mu}\bar b_\mu$, so comparing coefficients of $P_{\mu'}(x;t^{-1})$ gives the first formula. Expanding $e_\lambda=\sum_\nu K_{\nu'\lambda}s_\nu$ and $s_\nu=\sum_\kappa K_{\nu\kappa}(u)P_\kappa(x;u)$ \cite[III (2.6)]{Macdonald95} at $u=t^{-1}$ gives the second. The $K_{\nu\kappa}(t)$ have nonnegative coefficients \cite[III (6.5)]{Macdonald95}, and $d_{\lambda\mu}\in\mathbb{Q}[t]$ by \cref{thm:DS}, so $d_{\lambda\mu}\in\mathbb{N}[t]$. At $t=1$, $P_\kappa(x;1)=m_\kappa$ \cite[III \S2]{Macdonald95}, so $d_{\lambda\mu}(1)=[m_{\mu'}]e_\lambda=M_{\lambda\mu'}$. \qed

\begin{remark}[Prior work; no novelty claim for the matrix]\label{rem:H}
The matrix $d_{\lambda\mu}(t)$ is the classical $e\to$Hall--Littlewood transition matrix: writing $e_\lambda=\sum_\kappa M(e,P)_{\lambda\kappa}P_\kappa(x;t)$, \cref{cor:dmatrix} says $d_{\lambda\mu}(t)=t^{\,n(\mu')-n(\lambda')}M(e,P)_{\lambda\mu'}(t^{-1})$, and its positivity and the value at $t=1$ are in \cite[\S3.2]{Kirillov98}. What is new is only the identification of this matrix with the $s$-leading coefficients of $\star$, i.e.\ \cref{thm:H}. Hikita's \cite[Lemma~6.3]{Hikita25} treats the unrescaled limit $q\to\infty$.
\end{remark}

\begin{remark}[Other edges]\label{rem:edges}
The edges $t\to\infty$, $s\to\infty$ and $t=0$ of the $(s,t)$-square, and the reflection $(s,t)\mapsto(1/s,1/t)$, can be treated through a conjugation of $e_k\star(\cdot)$ by an operator that is diagonal on Macdonald polynomials; such a conjugation is implicit in \cite{DFK19}. We do not develop this here and use none of it: the only other edge that enters below is $t=0$, and for it we cite \cite{DFK18} (\cref{rem:t0input}).
\end{remark}


%% ===================================================================
\section{The block valuation law at $s=1$}\label{sec:blocklaw}
%% ===================================================================

\subsection{Taylor pieces}
We write $\val$ for the $(s-1)$-adic valuation on $\mathbb{Q}(t)[s]$, $I=(e_1,e_2,\dots)$ for the augmentation ideal (so $I^b$ is spanned by the $e_\mu$ with $\ell(\mu)\ge b$), and $\lin G:=[e_n]G$ for $G$ homogeneous of degree $n$; $\lin$ vanishes on $I^2$. For $\rho$ a partition, a homogeneous symmetric function is \emph{of type $\rho$} if it lies in $\mathrm{span}\{e_\nu:\nu\trianglerighteq\rho\}$; by \cref{lem:e-m,lem:Dj} this means $D_j(P)\le\sum_i\min(\rho_i,j)$ for all $j$. For a set $C$ of positions, $\lambda_C$ is the partition formed by the parts $\lambda_i$, $i\in C$.

In \cref{thm:subset} the variable $s$ enters only through $F|_{x_i\mapsto sx_i\,(i\in A)}=s^{\Delta_A}F$, $\Delta_A=\sum_{i\in A}x_i\partial_i$. Since $s^{\Delta_A}x^\alpha=s^{\alpha\cdot1_A}x^\alpha=\sum_p(s-1)^p\binom{\Delta_A}{p}x^\alpha$,
\begin{equation}\label{eq:taylor}
E_k=\sum_{p\ge0}(s-1)^pE_k^{(p)},\qquad E_k^{(p)}=\sum_{|A|=k}c_Ax_A\binom{\Delta_A}{p}.
\end{equation}
As $E_kF$ is a symmetric polynomial in $x$ and $s$ (\cref{thm:DS}), each $E_k^{(p)}$ preserves $\Lam_m$. Since $\binom{\Delta_A}p$ is a polynomial of degree $p$ in a derivation, and multiplication by $c_Ax_A$ preserves this, $E_k^{(p)}$ is a differential operator of order $\le p$: all $(p+1)$-fold commutators $[\cdots[E_k^{(p)},a_0],\dots,a_p]$ with multiplication operators vanish. Moreover $E_k^{(0)}$ is multiplication by $e_k$ (by the proof of \cref{thm:DS}, $(E_kF)|_{s=1}=e_kF|_{s=1}$), and $E_k^{(p)}(1)=0$ for $p\ge1$.

\subsection{First order}
\begin{theorem}[First order is a biderivation]\label{thm:bider}
Let $D_k:=E_k^{(1)}$ and $B(f,g):=\partial_s(f\star g)|_{s=1}$ for $f,g\in\Lam_m$ with coefficients in $\mathbb{Q}(t)$ (independent of $s$). Then $D_k$ is a derivation of $\Lam_m$, and
\[ B(f,g)=\sum_{k,l}M_{kl}\,\frac{\partial f}{\partial e_k}\frac{\partial g}{\partial e_l},\qquad M_{kl}:=D_k(e_l)=B(e_k,e_l)=M_{lk}. \]
So $B$ is a symmetric biderivation, the carr\'e du champ of $L=\frac12\sum M_{kl}\partial^2/\partial e_k\partial e_l$.
\end{theorem}
\begin{proof}
$D_k=\sum_Ac_Ax_A\Delta_A$ is a sum of derivations. By \cref{thm:DS} the $e^\star_\rho$ form a basis whose transition matrix to the $e_\rho$ is unitriangular at $s=1$; write $f=\sum_\rho a_\rho(s)e^\star_\rho$ with $a_\rho$ regular at $s=1$. By associativity, $f\star g=\sum_\rho a_\rho(s)E_{\rho_1}\cdots E_{\rho_\ell}g$. By \eqref{eq:taylor},
\[ B(f,g)=\sum_\rho a'_\rho(1)e_\rho g+\sum_\rho a_\rho(1)\sum_i\Bigl(\prod_{j\ne i}e_{\rho_j}\Bigr)D_{\rho_i}g . \]
Applying the same computation to $f=\sum_\rho a_\rho(s)E_{\rho_1}\cdots E_{\rho_\ell}(1)$ itself, which does not depend on $s$, and using $D_k(1)=0$, gives $\sum_\rho a'_\rho(1)e_\rho=0$. Since $f|_{s=1}=\sum_\rho a_\rho(1)e_\rho$, the remaining sum is $\sum_k\frac{\partial f}{\partial e_k}D_k(g)$, and $D_k(g)=\sum_l\frac{\partial g}{\partial e_l}D_k(e_l)$. Symmetry of $M$ is commutativity of $\star$ \cite{Hikita25}.
\end{proof}

\begin{proposition}[The matrix $M$]\label{prop:M}
For $k\ge r\ge0$, with $L(a,b):=\dfrac{(1-t^{a+b})(t^{ab}-1)}{(1-t^a)(1-t^b)}$,
\[ M_{kr}=B(e_k,e_r)=r\,e_ke_r+\sum_{j=1}^{r}L(k-r+j,\,j)\,e_{k+j}e_{r-j}. \]
Consequently $[(s-1)^1]e^\star_\lambda=\sum_{i<j}M_{\lambda_i\lambda_j}\,e_{\lambda\setminus\{\lambda_i,\lambda_j\}}$.
\end{proposition}
\begin{proof}
$M_{kr}=[(s-1)^1]\,e_k\star e_r$. Use the form of the Pieri rule in \cref{cor:pieri} (\cref{sec:box}; its proof does not use \cref{prop:M}): $e_k\star e_r=\sum_{y=0}^{r}s^y\mathcal C_{k-y,r-y}e_{k+r-y}e_y$. Since $\pi_a(b)|_{s=1}=\delta_{b0}$, every $\mathcal C_{a,b}$ with $b\ge1$ vanishes at $s=1$; so the $y=r$ term $s^re_ke_r$ contributes $r\,e_ke_r$, and the term $y=r-j$, $j\ge1$, contributes $\partial_s\mathcal C_{a,j}|_{s=1}$ with $a=k-r+j$. Now $\partial_s\pi_a(j)|_{s=1}=[u^j]\sum_{m<a}\frac{ut^m}{1+ut^m}=(-1)^{j-1}[a]_{t^j}$, so $\partial_s\mathcal C_{a,j}|_{s=1}=-\frac{[a+j]}{[a]}[a]_{t^j}=L(a,j)$. For the second statement, $[(s-1)^1]E_{\lambda_1}\cdots E_{\lambda_\ell}(1)=\sum_i\prod_{j<i}e_{\lambda_j}\,D_{\lambda_i}\bigl(\prod_{j>i}e_{\lambda_j}\bigr)$; expand with the Leibniz rule.
\end{proof}

\subsection{Blocks}
\begin{lemma}[Polarization]\label{lem:polar}
Let $D$ be a linear operator on a commutative ring and $g_1,\dots,g_r$ elements of it. For $S\subseteq[r]$ let $\delta_S:=[\cdots[D,g_{s_1}],\dots,g_{s_q}](1)=\sum_{T\subseteq S}(-1)^{|S\setminus T|}\prod_{i\in S\setminus T}g_i\,D\bigl(\prod_{i\in T}g_i\bigr)$. Then
\[ D(g_1\cdots g_r)=\sum_{S\subseteq[r]}\Bigl(\prod_{i\notin S}g_i\Bigr)\delta_S , \]
and if $D$ has order $\le p$, then $\delta_S=0$ for $|S|>p$, and for $|S|=p$ the map $(g_i)_{i\in S}\mapsto\delta_S$ is a derivation in each argument.
\end{lemma}
\begin{proof}
Substituting the second expression for $\delta_S$, the right side is $\sum_T\prod_{i\notin T}g_i\,D(\prod_{i\in T}g_i)\sum_{S\supseteq T}(-1)^{|S\setminus T|}$, and the inner sum is $[T=[r]]$. Vanishing is the definition of order $\le p$. For $|S|=p$, $D':=[\cdots[D,g_{s_2}],\dots,g_{s_p}]$ has order $\le1$, so $D'(ab)=aD'(b)+bD'(a)-abD'(1)$, i.e.\ $[D',ab](1)=a[D',b](1)+b[D',a](1)$; and the iterated commutator is symmetric in its arguments.
\end{proof}

\begin{proposition}[Block expansion]\label{prop:block}
Let $\lambda$ have $\ell$ parts and $p_1,\dots,p_\ell\ge0$ with $\sum_ip_i=m$. Then
\[ E^{(p_1)}_{\lambda_1}\cdots E^{(p_\ell)}_{\lambda_\ell}(1)=\sum_\pi\prod_{C\in\pi}g^\pi_C , \]
a finite sum over set partitions $\pi$ of $[\ell]$ with at least $\ell-m$ blocks, where each $g^\pi_C$ is homogeneous of type $\lambda_C$.
\end{proposition}
\begin{proof}
Induction on the number of operators applied, from the right; the empty product $1$ has no blocks. Apply $D=E^{(p)}_k$, $k=\lambda_j$, $p=p_j$, to a term $\prod_{C\in\pi}g_C$. By \cref{lem:polar} the result is $\sum_{S\subseteq\pi,|S|\le p}\prod_{C\notin S}g_C\cdot\delta_S$, and we let $\pi'=(\pi\setminus S)\cup\{C'\}$ with $C'=\{j\}\cup\bigcup S$. Then $|\pi'|\ge|\pi|+1-p$, which gives the block count. Each $\delta_S$ is a combination of products of the $g_C$ and of $D$ applied to such products. Rescaling monomials does not enlarge monomial supports, so the degree step \cref{lem:degstep} holds for $E^{(p)}_k$ as well, and $D_j$ is additive on products; hence $D_j(\delta_S)\le\sum_{C\in S}\sum_{i\in C}\min(\lambda_i,j)+\min(\lambda_j,j)$, i.e.\ $\delta_S$ has type $\lambda_{C'}$.
\end{proof}

\begin{definition}
For $\lambda,\mu\vdash n$, $\kappa(\lambda,\mu)$ is the largest $r$ such that $\lambda=\lambda^1\sqcup\cdots\sqcup\lambda^r$ and $\mu=\mu^1\sqcup\cdots\sqcup\mu^r$ as multisets of parts, with every block nonempty and $\mu^i\trianglerighteq\lambda^i$ (so $|\mu^i|=|\lambda^i|$); $\kappa=-\infty$ if $\mu\not\trianglerighteq\lambda$. $\mu$ is a \emph{coarsening} of $\lambda$ if $\kappa(\lambda,\mu)=\ell(\mu)$.
\end{definition}

\subsection{The block valuation law}
\begin{theorem}[Block valuation law]\label{thm:blocklaw}
For all $\lambda,\mu\vdash n$ and all $t$, $\val(c_{\lambda\mu})\ge\ell(\lambda)-\kappa(\lambda,\mu)$. For $\mu\trianglerighteq\lambda$, $\mu\neq\lambda$, equality holds over $\mathbb{Q}(t)$ and at $t=0$:
\[ \val(c_{\lambda\mu})=\ell(\lambda)-\kappa(\lambda,\mu). \]
At $t=0$ the leading coefficient is $(-1)^{\ell-\kappa}N(\lambda,\mu)$ with $N(\lambda,\mu)\ge1$ a count of minimal raising-operator configurations.
\end{theorem}

\begin{proof}[Proof of the lower bound]
By \eqref{eq:taylor}, $[(s-1)^m]e^\star_\lambda=\sum_{\sum p_i=m}E^{(p_1)}_{\lambda_1}\cdots E^{(p_\ell)}_{\lambda_\ell}(1)$. By \cref{prop:block} each term is a sum of products $\prod_{C\in\pi}g_C$ with $|\pi|\ge\ell-m$ and $g_C$ of type $\lambda_C$, and such a product lies in $\mathrm{span}\{e_{\sqcup_C\nu^C}:\nu^C\trianglerighteq\lambda_C\}$. So $e_\mu$ can occur only if $\kappa(\lambda,\mu)\ge|\pi|\ge\ell-m$. The argument works over $\mathbb{Q}(t)$ and for every specialisation $t\in\mathbb{Q}$.
\end{proof}

\begin{remark}[Input at $t=0$]\label{rem:t0input}
The $t=0$ edge $e^\star_\mu|_{t=0}=s^{n(\mu)}P_{\mu'}(x;1/s,0)=\omega\widetilde H_\mu(x;s)$, where $\widetilde H_\mu(x;s)=s^{n(\mu)}Q'_\mu(x;s^{-1})$, is \textbf{not ours}: it is assembled from \cite{DFK18} eq.~(5.15) (with $M_{k,1}=E_k|_{t=0}$), Cor.~5.8 in its level-one form (5.25), Cor.~5.18 (eq.~(5.27)) (numbering of arXiv:1505.01657v2), and \cite[VI (5.1), (4.14)(iv)]{Macdonald95}. In particular \cref{thm:blocklaw} does not depend on Cherednik's nonsymmetric theory.
\end{remark}

\begin{proof}[Proof of equality]
By \cref{rem:t0input}, $c_{\lambda\mu}(s,0)=[h_\mu]\,s^{n(\lambda)}Q'_\lambda(x;s^{-1})$. By \cite[III (2.15$'$) with \S5 Ex.~7(a)]{Macdonald95}, $Q'_\lambda(x;u)=\prod_{i<j\le\ell}\frac{1-R_{ij}}{1-uR_{ij}}h_\lambda$, where $R_{ij}h_\alpha=h_{\alpha+e_i-e_j}$, $h_\alpha=h_{\alpha_1}h_{\alpha_2}\cdots$ and $h_r=0$ for $r<0$. Expand $\frac{1-R}{1-uR}=1+\sum_{a\ge1}(u-1)u^{a-1}R^a$; at $u=s^{-1}$ the coefficient of $R^a$ is $-(s-1)s^{-a}$. A \emph{configuration} is a function $\mathbf m$ on pairs $i<j\le\ell$ with values in $\mathbb{Z}_{\ge0}$; its edge set is $E(\mathbf m)=\{ij:\mathbf m_{ij}\ge1\}$ and its result is $\alpha(\mathbf m)=\lambda+\sum\mathbf m_{ij}(e_i-e_j)$. Then
\[ c_{\lambda\mu}(s,0)=\sum_{\mathbf m:\ \alpha(\mathbf m)\ge0,\ \mathrm{sort}\,\alpha(\mathbf m)=\mu}(-1)^{|E(\mathbf m)|}(s-1)^{|E(\mathbf m)|}s^{a(\mathbf m)},\qquad a(\mathbf m)\in\mathbb{Z}, \]
a finite sum (the conditions $\alpha_\ell\ge0$, $\alpha_{\ell-1}\ge0,\dots$ bound the $\mathbf m_{ij}$ successively). Every configuration contributes $(-1)^{|E|}$ at order $(s-1)^{|E|}$, so there is no cancellation at the lowest order: $\val(c_{\lambda\mu}(s,0))=E_{\min}$, the least $|E(\mathbf m)|$ over admissible $\mathbf m$, with leading coefficient $(-1)^{E_{\min}}N(\lambda,\mu)$, $N$ the number of minimisers.

$E_{\min}\ge\ell-\kappa$: let $C$ be a connected component of the graph $([\ell],E(\mathbf m))$. Transfers stay inside $C$ and move mass from a larger index to a smaller one, and $\lambda$ is weakly decreasing; so the partial sums of $\alpha|_C$ in index order dominate those of $\lambda_C$, and sorting only increases them. Hence the components form a compatible block decomposition, there are at most $\kappa$ of them, and $|E|\ge\ell-\kappa$.

$E_{\min}\le\ell-\kappa$: take a decomposition with $\kappa$ blocks. By maximality each block $(\rho,\nu)=(\lambda_C,\mu^C)$ has no compatible refinement. Let $c_1<\dots<c_L$ be the positions of $C$, pad $\nu$ with zeros to length $L$, and put $B_r=\nu_1+\dots+\nu_r$, $P_r=\rho_1+\dots+\rho_r$. Then $B_r>P_r$ for $1\le r<L$: $\ge$ holds by dominance, and $B_r=P_r$ would split $(\rho,\nu)$ into the compatible pieces $(\rho_{\le r},\nu_{\le r})$ and $(\rho_{>r},\nu_{>r})$, both of positive size. Put $\mathbf m_{c_rc_{r+1}}=B_r-P_r\ge1$. Position $c_r$ then ends with $\rho_r+(B_r-P_r)-(B_{r-1}-P_{r-1})=\nu_r\ge0$. This uses $L-1$ edges per block, $\ell-\kappa$ in all.

Finally, $c_{\lambda\mu}\in\mathbb{Q}[s,t]$ (\cref{thm:DS}); write $c_{\lambda\mu}=\sum_j(s-1)^ja_j(t)$. The lower bound gives $a_j=0$ for $j<\ell-\kappa$, and $a_{\ell-\kappa}(0)=(-1)^{\ell-\kappa}N(\lambda,\mu)\ne0$, so $a_{\ell-\kappa}\ne0$.
\end{proof}

\begin{remark}[Novelty]
We have not found the block law stated anywhere. Its $t=0$ case is classical in substance: the raising-operator formula follows from \cite[eq.~(11)]{DLT94} and Jacobi--Trudi, after which the valuation count above is a short exercise. What is new is the lower bound for all $t$ (from the subset formula alone), the exactness over $\mathbb{Q}(t)$, and the merge weights below.
\end{remark}

\subsection{Coarsenings and merge weights}
\begin{theorem}[Coarsenings: merge-history expansion]\label{thm:coarse}
Let $\mu\neq\lambda$ be a coarsening of $\lambda$ and $m=\ell(\lambda)-\ell(\mu)$. Then $\val(c_{\lambda\mu})=m$, and $\Lead_{\lambda\mu}(t):=[(s-1)^m]c_{\lambda\mu}$ equals
\[ \Lead_{\lambda\mu}(t)=\sum_{\text{tight histories }H\text{ ending at }\mu}\ \prod_{\text{merges in }H}W_k(J),\qquad \Lead_{\lambda\mu}(0)=(-1)^m\#\{H\}\neq0, \]
where a (tight) history processes $\lambda_\ell,\dots,\lambda_1$, each part $k$ either opening a block of size $k$ or merging $p\ge1$ existing blocks of sizes $J=(j_1,\dots,j_p)$ into one block of size $k+|J|$ via the order-$p$ piece $E_k^{(p)}$, it is tight if the orders add up to $m$, it ends at $\mu$ if the final block sizes are $\mu$, and $W_k(J):=\lin\,[\cdots[E_k^{(p)},e_{j_1}],\dots,e_{j_p}](1)$.
\end{theorem}
\begin{proof}
Since $\ell(\mu)=\ell-m$, the coefficient of $e_\mu$ in $[(s-1)^m]e^\star_\lambda$ can be read modulo $I^{\ell-m+1}$. In \cref{prop:block}, a product over $|\pi|>\ell-m$ blocks of positive-degree factors lies in $I^{\ell-m+1}$; so only the terms where every step merges exactly $p_j$ blocks survive, i.e.\ tight histories. At a tight step, $\delta_S$ is a derivation in each argument (\cref{lem:polar}), so $\delta_S(bc,\dots)\in I^2$ for $b,c\in I$, and modulo $I^2$, $\delta_S(g_1,\dots,g_p)$ depends only on the linear parts $g_i\equiv\gamma_ie_{n_i}$: it is $\prod_i\gamma_i\,W_k(n_1,\dots,n_p)\,e_n$. An opening step contributes $e_k$, with $\gamma=1$. The final product of blocks is $\prod_C\gamma_Ce_{|C|}$ modulo $I^{\ell-m+1}$, which gives the formula. By \cref{thm:W} below, $W_k(J)|_{t=0}=(-1)^{|J|}$, so every tight history contributes $(-1)^m$ at $t=0$. Histories exist: in each block of the coarsening, let the last-processed part merge all the others, opened earlier as singletons. So $\Lead_{\lambda\mu}(0)\ne0$, and $\val(c_{\lambda\mu})=m$ by the lower bound in \cref{thm:blocklaw}.
\end{proof}

Let $T_kf:=\sum_{|A|=k}c_Ax_Af(x_A)$ for $f\in\Lam_k$, where $f(x_A)$ is $f$ evaluated at the $k$ variables $x_a$, $a\in A$. (This is unrelated to the Hecke operators of \cref{sec:hikita}.)

\begin{theorem}[Merge weights]\label{thm:W}
With $n=k+|J|$ and $p=|J|$,
\[ W_k(J)=(-1)^{p}\,\frac{[n]_t}{[k]_t}\prod_{j\in J}[k]_{t^{j}}. \]
Consequently $(-1)^m\Lead_{\lambda\mu}(t)>0$ for coarsenings and every real $t>0$.
\end{theorem}
\begin{lemma}\label{lem:linT}
For $f\in\Lam_k$ homogeneous of degree $d$ and $n=k+d$: \[\lin T_k(f)=(-1)^d\frac{[n]_t}{[k]_t}f(1,t,\dots,t^{k-1}).\]
\end{lemma}
\begin{remark}[Novelty]
\Cref{lem:linT} is a short corollary of \cite[III.7 Ex.~2, III (2.2)]{Macdonald95} and is probably known.
\end{remark}

\begin{lemma}\label{lem:parabolic}
For $\ell(\rho)\le k\le m$, $T_k\bigl(P_\rho(x_1,\dots,x_k;t)\bigr)=P_{\rho+1^k}(x_1,\dots,x_m;t)$.
\end{lemma}
\begin{proof}
By \cite[III (2.2)]{Macdonald95}, $P_\lambda=\sum_{w\in S_m/S_m^\lambda}w\bigl(x^\lambda\prod_{i<j,\,\lambda_i>\lambda_j}\frac{x_i-tx_j}{x_i-x_j}\bigr)$, $S_m^\lambda$ the stabiliser of $\lambda$. For $\lambda=\rho+1^k$ the positive entries are the first $k$; a coset is an image set $A=w([k])$ together with a coset of $S_k/S_k^\rho$; the pairs with $\lambda_i>\lambda_j$ are those with $i\le k<j$, giving $c_A$, and those with $i,j\le k$, $\rho_i>\rho_j$; and $x^\lambda=x_{[k]}x_{[k]}^\rho$. Grouping by $A$, the inner sum is \cite[III (2.2)]{Macdonald95} for $P_\rho$ in the variables $x_A$.
\end{proof}

\begin{lemma}\label{lem:linP}
For $\lambda\vdash n$ with $\ell(\lambda)=k$, $\lin P_\lambda(x;t)=(-1)^{n-k}\frac{[n]_t}{[k]_t}P_{\lambda-1^k}(1,t,\dots,t^{k-1};t)$.
\end{lemma}
\begin{proof}
Since $[p_n]e_n=(-1)^{n-1}/n$ and $e_\nu$ with $\ell(\nu)\ge2$ has no $p_n$ term, $\lin G=(-1)^{n-1}\langle G,p_n\rangle=(-1)^{n-1}(1-t^n)\langle G,p_n\rangle_t$. By $\langle P_\lambda,Q_\mu\rangle_t=\delta_{\lambda\mu}$ and $Q_\lambda=b_\lambda P_\lambda$ \cite[III (4.9)]{Macdonald95}, $\langle P_\lambda,p_n\rangle_t=[P_\lambda]p_n/b_\lambda(t)$, and $[P_\lambda]p_n=t^{n(\lambda)}\varphi_{k-1}(t^{-1})=(-1)^{k-1}t^{n(\lambda)-\binom k2}\varphi_{k-1}(t)$ \cite[III.7 Ex.~2]{Macdonald95}, where $\varphi_r(t)=\prod_{a=1}^r(1-t^a)$. On the other side, with $\rho=\lambda-1^k$, \cite[III.2 Ex.~1]{Macdonald95} gives $P_\rho(1,\dots,t^{k-1})=t^{n(\rho)}\varphi_k/\prod_{i\ge0}\varphi_{m_i(\rho)}$ with $m_0(\rho)=k-\ell(\rho)$; here $m_i(\rho)=m_{i+1}(\lambda)$, so the denominator is $b_\lambda$, and $n(\rho)=n(\lambda)-\binom k2$. Since $\frac{[n]}{[k]}\varphi_k=(1-t^n)\varphi_{k-1}$, the two sides agree.
\end{proof}

\begin{proof}[Proof of \cref{lem:linT}]
The $P_\rho(x_1,\dots,x_k;t)$ with $\ell(\rho)\le k$ form a basis of $\Lam_k$ \cite[III (2.7)]{Macdonald95}; for $f=P_\rho$, $\rho\vdash d$, combine \cref{lem:parabolic,lem:linP}.
\end{proof}

\begin{proof}[Proof of \cref{thm:W}]
Let $J=(j_1,\dots,j_p)$, $d=|J|$. First, $W_k(J)=\lin E_k^{(p)}(e_{j_1}\cdots e_{j_p})$: in \cref{lem:polar}, $\delta_\varnothing=E^{(p)}_k(1)=0$, and for $\varnothing\ne S\ne[p]$ the term $\prod_{i\notin S}e_{j_i}\,\delta_S$ lies in $I^2$. Next, with $E(u)=\sum_re_ru^r$ and $R(u):=\prod_{a\in A}\frac{1+sux_a}{1+ux_a}$ read as a function of $x_A$,
\[ E_k\Bigl(\prod_{i=1}^pE(u_i)\Bigr)=\prod_iE(u_i)\cdot T_k\Bigl(\prod_iR(u_i)\Bigr),\qquad \log R(u)=\sum_{r\ge1}(-1)^{r-1}\frac{s^r-1}{r}u^rp_r(x_A). \]
Take $[u_1^{j_1}\cdots u_p^{j_p}]$ and $[(s-1)^p]$. Terms in which some $e_{r_i}$, $r_i\ge1$, is taken from $\prod E(u_i)$ lie in $I^2$. In the remaining term, each $[u_i^{j_i}]R(u_i)$ has $(s-1)$-valuation $\ge1$, with $(s-1)^1$-coefficient $(-1)^{j_i-1}p_{j_i}(x_A)$. Hence $W_k(J)=(-1)^{d-p}\lin T_k(p_J)$, and \cref{lem:linT} with $p_j(1,\dots,t^{k-1})=[k]_{t^j}$ gives the formula. For $t>0$ every $W_k(J)$ has sign $(-1)^{|J|}$, and the orders along a tight history add up to $m$.
\end{proof}


%% ===================================================================
\section{Leading coefficients: graph cumulants and blocks}\label{sec:leads}
%% ===================================================================
Throughout this section $\mu$ is a coarsening of $\lambda$, $\ell=\ell(\lambda)$, $m=\ell-\ell(\mu)$, and the parts of $\lambda$ are indexed by positions $[\ell]$. Everything rests on \cref{thm:coarse,thm:W}.

\begin{lemma}[Histories are increasing forests]\label{lem:forest}
Tight histories of $\lambda$ are in bijection with increasing forests on $[\ell]$ (every child larger than its parent): when position $i$ merges existing blocks, $i$ becomes the parent of their roots. The trees are the final blocks, and the orders automatically add up to $\ell-\#\text{trees}$.
\end{lemma}
\begin{proof}
Every existing block at step $i$ was built from positions $>i$, so all edges increase. Conversely, an increasing forest determines the history: at step $i$, merge the subtrees rooted at the children of $i$, which are complete by then. Each edge accounts for one unit of order.
\end{proof}

\begin{lemma}[Telescoping]\label{lem:telescope}
For an increasing tree $T$ on a set $C$ of positions, let $N_v$ be the $\lambda$-mass of the subtree at $v$. The product of the merge weights along the history of $T$ is
\[ \prod_vW_{\lambda_v}\bigl((N_c)_{c\ \text{child of}\ v}\bigr)=\frac{1-t^{|\lambda_C|}}{\prod_{i\in C}(1-t^{\lambda_i})}\prod_{\text{edges }v\to c}\bigl(t^{\lambda_vN_c}-1\bigr). \]
\end{lemma}
\begin{proof}
By \cref{thm:W}, the merge at $v$ contributes $(-1)^{\#\text{children}}\frac{1-t^{N_v}}{1-t^{\lambda_v}}\prod_{c}\frac{1-t^{\lambda_vN_c}}{1-t^{N_c}}$ (a leaf contributes $1$). Every non-root $c$ is a child exactly once, so its factor $1-t^{N_c}$ cancels against its own numerator, leaving the root mass $|\lambda_C|$; and $(-1)^{\#\text{edges}}\prod(1-x)=\prod(x-1)$.
\end{proof}

\begin{lemma}[Tree--graph identity]\label{lem:treegraph}
For indeterminates $x_{ij}=x_{ji}$ ($i\ne j\in[\ell]$),
\[ \sum_{\text{increasing trees }T\text{ on }[\ell]}\ \prod_{\text{edges }v\to c}\Bigl(\prod_{j\in T_c}x_{vj}-1\Bigr)=\sum_{\text{connected graphs }H\text{ on }[\ell]}\ \prod_{e\in H}(x_e-1), \]
where $T_c$ is the vertex set of the subtree at $c$.
\end{lemma}
\begin{proof}
Map a connected graph $H$ to an increasing tree recursively: the root is $\min V$, the children are the minima of the components of $H-\min V$, and recurse on each component. Fix $T$ and sum over the $H$ mapping to it. At each vertex $v$ and child $c$, the edges from $v$ into $T_c$ form an arbitrary nonempty set $S\subseteq T_c$ (nonempty since $T_c$ must be attached to $v$; otherwise free, since they do not change the components inside $T_c$); there are no edges between different child sets; and the edges inside $T_c$ are accounted for recursively. So the fibre sum factors, and $\sum_{\varnothing\ne S\subseteq T_c}\prod_{j\in S}(x_{vj}-1)=\prod_{j\in T_c}x_{vj}-1$.
\end{proof}

\begin{theorem}[Full-merge lead]\label{thm:G}
For $\lambda\vdash n$ with $\ell=\ell(\lambda)\ge2$,
\[ \Lead_{\lambda,(n)}(t)=\frac{1-t^n}{\prod_i(1-t^{\lambda_i})}\,K_\lambda(t),\qquad K_\lambda(t):=\sum_{H\text{ connected on }[\ell]}\ \prod_{ij\in H}\bigl(t^{\lambda_i\lambda_j}-1\bigr). \]
$K_\lambda$ is the connected part (joint cumulant) of $t^{e_2(\lambda)}=\prod_{i<j}(1+(t^{\lambda_i\lambda_j}-1))$.
\end{theorem}
\begin{proof}
By \cref{thm:coarse} with $\mu=(n)$ and \cref{lem:forest}, $\Lead_{\lambda,(n)}$ is a sum over increasing trees on $[\ell]$; by \cref{lem:telescope} each tree contributes $\frac{1-t^n}{\prod_i(1-t^{\lambda_i})}\prod_{v\to c}(t^{\lambda_vN_c}-1)$; and \cref{lem:treegraph} with $x_{ij}=t^{\lambda_i\lambda_j}$, for which $\prod_{j\in T_c}x_{vj}=t^{\lambda_vN_c}$, converts the tree sum into $K_\lambda$. The cumulant statement is the exponential formula.
\end{proof}
\begin{corollary}\label{cor:G}
(a) $\lambda=1^n$: $\Lead=(-1)^{n-1}[n]_t\,I_n(t)$ with $I_n$ the tree inversion enumerator \cite{MallowsRiordan68}. (b) $t=0$: $\Lead=(-1)^{\ell-1}(\ell-1)!$. (c) $t\to1$: $\Lead\to(-1)^{\ell-1}n^{\ell-1}$ (weighted Cayley). (d) $\operatorname{sgn}\Lead=(-1)^{\ell-1}$ for all $t>0$.
\end{corollary}
\begin{proof}
(a) For $\lambda=1^n$ the prefactor is $[n]_t/(1-t)^{n-1}$ and $K_\lambda=\sum_{H\text{ conn.}}(t-1)^{|H|}=(t-1)^{n-1}I_n(t)$ \cite{MallowsRiordan68}. (b) At $t=0$ the prefactor is $1$, and in the tree form every edge factor is $-1$; there are $(\ell-1)!$ increasing trees. (c) Put $\varepsilon=t-1$. In $K_\lambda$ only spanning trees contribute to the lowest order $\varepsilon^{\ell-1}$, with $t^{\lambda_i\lambda_j}-1=\lambda_i\lambda_j\varepsilon+O(\varepsilon^2)$, and $\sum_T\prod_{ij\in T}\lambda_i\lambda_j=\sum_T\prod_i\lambda_i^{\deg_T(i)}=(\prod_i\lambda_i)\,n^{\ell-2}$ by the vertex-weighted Cayley formula (each $i$ occurs $\deg_T(i)-1$ times in the Pr\"ufer code). The prefactor is $(-n\varepsilon)/((-\varepsilon)^\ell\prod_i\lambda_i)\cdot(1+O(\varepsilon))$. (d) is the case $\mu=(n)$, $m=\ell-1$, of the sign statement in \cref{thm:W}: each merge contributes $(-1)^p$ times a positive quantity, $\sum p=m$ along every history, and histories exist.
\end{proof}
\begin{remark}[Prior art --- graph half]
The tree, connected-graph and cumulant forms of $K_\lambda$ are a special case of Do{\l}{\k{e}}ga \cite[Prop.~2.1 and (11), Lemma~2.3]{Dolega19} ($K_\lambda=(t-1)^{\ell-1}T(1,t)$ for the Tutte polynomial of the multigraph with $\lambda_i\lambda_j$ edges between $i$ and $j$), cf.\ also \cite{GesselSagan96,MallowsRiordan68}; step (3) is a Penrose-type tree--graph identity \cite{Penrose67}. Ours: the identification of the $\star$-lead with prefactor$\cdot K_\lambda$ (\cref{thm:coarse,thm:W,thm:G}) and \cref{thm:F,thm:blockmult}. Haglund--Tewari's single-row cumulant \cite[Def.~7.1, Thm.~7.3]{HaglundTewari26} (their $q$ is our $t$) is the same graph half: since $\langle\widetilde H_{(m)},e_m\rangle=q^{\binom m2}$, their cumulant $\kappa(\lambda)$ (we write $C_\lambda$) satisfies $\langle C_\lambda,e_n\rangle=t^{n(\lambda')}(t-1)^{1-\ell}K_\lambda(t)$.
\emph{Separator.} Full-merge leads $L(\lambda)$ of any such theory are only defined up to the gauge $L(\lambda)\mapsto L(\lambda)\,a^{\ell(\lambda)}g(|\lambda|)\prod_if(\lambda_i)$ (rescaling generators, degree-wise normalisation, rescaling the expansion parameter). The first gauge invariant appears at $n=4$: $J:=L(1^4)L(2,2)/L(2,1,1)^2$ (the exponents of $a$, $g(4)$, $f(1)$ and $f(2)$ all cancel). For $\star$, \cref{thm:G} gives $J=\frac{(t^2+1)(t^3+3t^2+6t+6)}{(t+1)(t^2+t+2)^2}$, and Haglund--Tewari's $\langle C_\lambda,e_n\rangle$ gives the same value. Do{\l}{\k{e}}ga's $\oplus$-cumulant of columns (lowest $(q-1)$-coefficient of $[g_n]\widetilde H_{\lambda'}$, $g_k=\widetilde H_{1^k}$) gives, by computation, $J\equiv 3/2=J_\star(0)$, whereas $J_\star(1)=1$, the Cayley value. So up to gauge Do{\l}{\k{e}}ga's $\oplus$ is $t$-independent and agrees with $\star$ only at $t=0$, where the leads are partition-lattice M\"obius values (\cref{cor:G}(b)), while Haglund--Tewari's is the same graph object without $\star$.
\end{remark}

\begin{theorem}[Coarsening leads factor]\label{thm:F}
For $\mu$ a coarsening of $\lambda$,
\[ \Lead_{\lambda\mu}=\sum_{\pi}\ \prod_{B\in\pi}\Lead_{\lambda_B,(|\lambda_B|)}, \]
summed over set partitions $\pi$ of the positions $[\ell]$ whose sorted block sums equal $\mu$ (singletons contribute $1$). There is no interleaving factor.
\end{theorem}
\begin{proof}
By \cref{lem:forest}, tight histories ending at $\mu$ are increasing forests whose trees have sorted masses $\mu$, i.e.\ a set partition $\pi$ with block sums $\mu$ together with an increasing tree on each block. By \cref{lem:telescope} the weight is a product over the trees, each depending only on its own parts; summing over the trees on each block gives $\Lead_{\lambda_B,(|\lambda_B|)}$ by \cref{thm:G} (and $1$ for a singleton).
\end{proof}

\begin{theorem}[Block multiplicativity]\label{thm:blockmult}
Let $\kappa=\kappa(\lambda,\mu)\ge1$, $\ell=\ell(\lambda)$, $m=\ell-\kappa$. Then
\[ [(s-1)^m]\,c_{\lambda\mu}=\sum_{\pi}\ \sum_{(\nu^C)_{C\in\pi}}\ \prod_{C\in\pi}[(s-1)^{|C|-1}]\,c_{\lambda_C,\nu^C}, \]
where $\pi$ runs over set partitions of $[\ell]$ into exactly $\kappa$ blocks and $(\nu^C)$ over tuples with $\bigsqcup_C\nu^C=\mu$ and $\nu^C\trianglerighteq\lambda_C$; every factor then has $\kappa(\lambda_C,\nu^C)=1$. With \cref{thm:blocklaw} this is $\Lead_{\lambda\mu}$.
\end{theorem}
\begin{proof}
Expand $[(s-1)^m]e^\star_\lambda=\sum_{\sum p_i=m}E^{(p_1)}_{\lambda_1}\cdots E^{(p_\ell)}_{\lambda_\ell}(1)$ step by step with \cref{lem:polar}, as in \cref{prop:block}: this is a sum over histories $H$ (step $i$ chooses a set $S_i$ of current blocks with $|S_i|\le p_i$), each giving $\prod_{C\in\pi(H)}g^H_C$ with $g^H_C$ of type $\lambda_C$, and $|\pi(H)|=\ell-\sum|S_i|\ge\ell-m=\kappa$, with equality iff $|S_i|=p_i$ for all $i$ (tight). A non-tight term has more than $\kappa$ blocks, so its $e_\mu$-coefficient vanishes by maximality of $\kappa$. In a tight history, step $i$ only acts, through $\delta_{S_i}$, on the blocks it merges; so $g^H_C$ is the function produced by the restricted history on $C$, which is a single-block tight history of the composition $\lambda_C$ (positions of $C$ in their order), since $\sum_{i\in C}|S_i|=|C|-1$. Conversely, single-block tight histories on the blocks of $\pi$ interleave uniquely, by position, into a global tight history. Hence the tight part is $\sum_{|\pi|=\kappa}\prod_{C\in\pi}\mathrm{Conn}(\lambda_C)$, where $\mathrm{Conn}(\lambda_C)$ is the sum over single-block tight histories of $\lambda_C$. Applying the same expansion to $\lambda_C$ with $m=|C|-1$ shows $[e_\nu]\mathrm{Conn}(\lambda_C)=[e_\nu][(s-1)^{|C|-1}]e^\star_{\lambda_C}$ whenever $\kappa(\lambda_C,\nu)=1$. Finally, $[e_\mu]\prod_C\mathrm{Conn}(\lambda_C)=\sum_{(\nu^C)}\prod_C[e_{\nu^C}]\mathrm{Conn}(\lambda_C)$ over tuples with $\bigsqcup_C\nu^C=\mu$; factors with $\nu^C\not\trianglerighteq\lambda_C$ vanish by type, and $\kappa(\lambda_C,\nu^C)\ge2$ would refine $\pi$ to more than $\kappa$ compatible blocks. Commutativity of $\star$ \cite{Hikita25} identifies the coefficients for the composition $\lambda_C$ with those for the sorted partition.
\end{proof}
\begin{remark}
\Cref{thm:F} is the case where each $\nu^C$ is a single part. The factorisation has the shape of a Mayer (cluster) expansion, i.e.\ moment--cumulant inversion of the graph identity in \cref{thm:G}; the $\star$-specific content is that the connected leads are those of \cref{thm:G}, with no interleaving correction.
\end{remark}


%% ===================================================================
\section{Box Complement and the linear coefficient}\label{sec:box}
%% ===================================================================
In this section zero parts are allowed: $\lambda$ is padded with zeros to a prescribed length, $E_0=\mathrm{id}$, and $c_{\lambda\mu}$ refers to the nonzero parts. By \cref{lem:stable}, in $N$ variables (for every $N$, including $N$ below the degree) $e^\star_\lambda(x_1,\dots,x_N)=\sum_\mu c_{\lambda\mu}e_\mu(x_1,\dots,x_N)$, where $e_\mu=0$ if $\mu_1>N$.

\begin{theorem}[Box Complement]\label{thm:box}
Let $\lambda,\mu$ have at most $\ell$ parts (padded with zeros) and all parts $\le N$; put $N^\ell-\lambda:=(N-\lambda_\ell,\dots,N-\lambda_1)$. Then
\[ c_{N^\ell-\lambda,\,N^\ell-\mu}(s,t)=s^{N\binom{\ell}{2}-(\ell-1)|\lambda|}\,c_{\lambda\mu}(s,t). \]
\end{theorem}
\begin{lemma}[Inversion]\label{lem:inversion}
In $N$ variables, let $G$ be homogeneous of degree $g$ and $F(x):=e_N(x)^m\,G(1/x)$. Then
\[ E_{N-k}F=s^{m(N-k)-g}\,e_N^{m+1}\,(E_kG)(1/x). \]
\end{lemma}
\begin{proof}
Write $y=1/x$ and index the terms of $E_{N-k}$ by the complement $A$ ($|A|=k$) of the $(N-k)$-subset $B$. Then $c_B(x)=\prod_{i\in B,j\in A}\frac{y_j-ty_i}{y_j-y_i}=c_A(y)$; $x_B=e_N(x)\,y_A$; $F|_{x_i\mapsto sx_i\,(i\in B)}=s^{m(N-k)}e_N(x)^m\,G|_{y_i\mapsto s^{-1}y_i\,(i\in B)}$; and by homogeneity $G|_{y_i\mapsto s^{-1}y_i\,(i\in B)}=s^{-g}\,G|_{y_a\mapsto sy_a\,(a\in A)}$. Summing over $A$ gives $s^{m(N-k)-g}e_N^{m+1}\sum_Ac_A(y)y_AG|_{y_a\mapsto sy_a}$.
\end{proof}

\begin{proof}[Proof of \cref{thm:box}]
Let $G_j:=E_{\lambda_{\ell-j+1}}\cdots E_{\lambda_\ell}(1)$, of degree $D_j=\lambda_\ell+\dots+\lambda_{\ell-j+1}$, and $F_j:=E_{N-\lambda_{\ell-j+1}}\cdots E_{N-\lambda_\ell}(1)$, in $N$ variables. By \cref{lem:inversion} and induction, $F_j=s^{\sigma_j}e_N^jG_j(1/x)$ with $\sigma_0=0$ and $\sigma_{j+1}=\sigma_j+j(N-\lambda_{\ell-j})-D_j$. Writing $p_i=\lambda_{\ell+1-i}$,
\[ \sigma_\ell=N\binom\ell2-\sum_i(i-1)p_i-\sum_i(\ell-i)p_i=N\binom\ell2-(\ell-1)|\lambda| . \]
By commutativity of $\star$, $F_\ell=e^\star_{N^\ell-\lambda}$. Since $e_{N-r}(x)=e_N(x)e_r(1/x)$ in $N$ variables, $e_N^\ell e_\mu(1/x)=e_{N^\ell-\mu}(x)$ for $\ell(\mu)\le\ell$, and every $\mu$ in the support of $e^\star_\lambda$ has $\ell(\mu)\le\ell(\lambda)\le\ell$ (\cref{thm:DS}). Hence $e^\star_{N^\ell-\lambda}=s^{\sigma_\ell}\sum_{\mu_1\le N}c_{\lambda\mu}\,e_{N^\ell-\mu}(x)$. The map $\mu\mapsto N^\ell-\mu$ is a bijection between partitions with at most $\ell$ parts, all $\le N$; every $\nu$ in the support of $e^\star_{N^\ell-\lambda}$ has at most $\ell$ parts; and the $e_\nu$ with $\nu_1\le N$ form a basis of $\Lam_N$. Compare coefficients.
\end{proof}
\begin{corollary}[Column Lemma; lead invariance]\label{cor:column}
$c_{\lambda+1^\ell,\,\mu+1^\ell}=s^{\binom{\ell}{2}}c_{\lambda\mu}$. Hence $\val$ and $\Lead$ are invariant under $(\lambda,\mu)\mapsto(\lambda+1^\ell,\mu+1^\ell)$ and $(\lambda,\mu)\mapsto(N^\ell-\lambda,N^\ell-\mu)$; e.g.\ $(2,2,2)\to(3,3)$ and $(2,2,2)\to(4,1,1)$ have the lead $[3]_t(t+2)$ of $(1,1,1)\to(3)$.
\end{corollary}
\begin{proof}
Complement with $N$ and then with $N+1$: $(N+1)^\ell-(N^\ell-\lambda)=\lambda+1^\ell$, and the exponents add up to $\bigl(N\binom\ell2-(\ell-1)|\lambda|\bigr)+\bigl((N+1)\binom\ell2-(\ell-1)(N\ell-|\lambda|)\bigr)=\binom\ell2$. Since $s^\sigma$ is a unit of $\mathbb{Q}[[s-1]]$ with constant term $1$, valuations and leading coefficients are preserved. For the example, $(2,2,2)\to(4,1,1)$ is $(1,1,1)\to(3,0,0)$ plus a column, and $(2,2,2)\to(3,3)$ is the complement with $N=3$ of $(1,1,1)\to(3,0,0)$; the lead of $(1,1,1)\to(3)$ is $[3]_t(t+2)$ by \cref{thm:G}.
\end{proof}
\begin{remark}[Prior art]
The inversion lemma follows in a few lines from \cite[Rem.~3.3]{DFK19} ($S\Gamma_IS=\Gamma_I^{-1}$ under $x\mapsto x^{-1}$) plus subset complement; stability is \cite[Prop.~3.8, Cor.~3.10]{Hikita25}. We have not found \cref{thm:box} or \cref{cor:column} for $\star$ in the literature.
\end{remark}

\begin{theorem}[Plethystic linear coefficient]\label{thm:lin}
For $G$ homogeneous of degree $d$: $\lin(e_k\star G)=(-1)^d\frac{[k+d]_t}{[k]_t}\,G[(s-1)[k]_t]$, where $G\mapsto G[(s-1)[k]_t]$ is the ring map $p_r\mapsto(s^r-1)[k]_{t^r}$.
\end{theorem}
\begin{proof}
By linearity take $G=e_{j_1}\cdots e_{j_p}$, $d=\sum j_i$. As in the proof of \cref{thm:W}, $E_k(e_J)=\sum_{w\le J}e_{J-w}\,T_k\bigl([u^w]\prod_iR(u_i)\bigr)$. For $w\ne J$ the term is a product of an element of $I$ and a symmetric function of positive degree, hence lies in $I^2$. For $w=J$, \cref{lem:linT} applied to the symmetric function $f=\prod_i[u_i^{j_i}]R(u_i)$ of $x_A$ gives $(-1)^d\frac{[k+d]}{[k]}f(1,t,\dots,t^{k-1})$, and $[u^j]R(u)$ at $x_A=(1,t,\dots,t^{k-1})$ is $\pi_k(j):=[u^j]\prod_{m=0}^{k-1}\frac{1+sut^m}{1+ut^m}$. Taking logarithms, $e_j\mapsto\pi_k(j)$ is the ring map $p_r\mapsto(s^r-1)[k]_{t^r}$.
\end{proof}
\begin{corollary}[Pieri rule, all $s$]\label{cor:pieri}
$e_k\star e_r=\sum_{y=0}^{\min(k,r)}s^y\mathcal C_{k-y,r-y}\,e_{k+r-y}e_y$ with $\mathcal C_{a,0}=\mathcal C_{0,b}=1$, $\mathcal C_{a,b}=(-1)^b\frac{[a+b]}{[a]}\pi_a(b)$, $\pi_a(j)=[u^j]\prod_{m=0}^{a-1}\frac{1+sut^m}{1+ut^m}$.
\end{corollary}
\begin{proof}
By \cref{thm:DS} and commutativity, the support of $e_k\star e_r$ consists of $e_{k+r-y}e_y$ with $y\le\min(k,r)$. Applying \cref{cor:column} with $\ell=2$ $y$ times, $c_{(k,r),(k+r-y,y)}=s^y\,c_{(k-y,r-y),(k+r-2y,0)}=s^y\lin\bigl(e_{k-y}\star e_{r-y}\bigr)$. If $k-y=0$ this is $s^y$; otherwise \cref{thm:lin} with $G=e_{r-y}$ gives $s^y(-1)^{r-y}\frac{[k+r-2y]}{[k-y]}\pi_{k-y}(r-y)=s^y\mathcal C_{k-y,r-y}$.
\end{proof}
\begin{remark}
This proof of the Pieri rule is independent of the one in \cref{sec:pieri}: it uses neither the parabolic kernel computation of \cref{prop:Ak} beyond the subset formula, nor the recursion of \cref{thm:ellcol}. For $r\ge k$ the monomials $e_be_{r+k-b}$, $0\le b\le k$, are distinct, so comparing with \cref{cor:pieri-ekr} gives $F_n(t^m)=(-1)^m\frac{[n+m]}{[n]}\pi_n(m)$ for $1\le n\le m$. Hikita's \cite[Thm.~3.12]{Hikita25} is the case $k=1$.
\end{remark}


%% ===================================================================
\section{Second order: every valuation-$2$ lead}\label{sec:v2}
%% ===================================================================

By \cref{thm:blockmult}, valuation-$2$ leads reduce to connected leads with $\ell(\lambda)=3$, $\kappa=1$; by \cref{cor:column} to two-part $\mu$. With $T_a$ as in \cref{lem:linT}, put
\[ \Gamma_a(e_b,e_c):=E_a^{(2)}(e_be_c)-e_bE_a^{(2)}(e_c)-e_cE_a^{(2)}(e_b)=\sum_{1\le r\le b,\ 1\le q\le c}(-1)^{r+q}e_{b-r}e_{c-q}\,T_a(p_rp_q). \]
Indeed, with $R(u)$ as in the proof of \cref{thm:W}, $R(u)=\prod_{x\in A}\bigl(1+(s-1)y_x(u)\bigr)$ with $y_x(u)=\frac{ux_x}{1+ux_x}$, so $[(s-1)^2]R(u)R(v)=e_2(Y_u)+e_2(Y_v)+p_1(Y_u)p_1(Y_v)$ in the alphabets $Y_u=\{y_x(u)\}_{x\in A}$, $Y_v$; hence $E_a^{(2)}(E(u)E(v))-E(u)E_a^{(2)}E(v)-E(v)E_a^{(2)}E(u)=E(u)E(v)\,T_a\bigl(p_1(Y_u)p_1(Y_v)\bigr)$, and $p_1(Y_u)=\sum_{r\ge1}(-1)^{r-1}u^rp_r(x_A)$. Take $[u^bv^c]$.

\begin{proposition}\label{prop:reduce}
If $(a,b,c)$ is any ordering of $\lambda$ and $\kappa(\lambda,\mu)=1$, then $[(s-1)^2]c_{\lambda\mu}=[e_\mu]\bigl(\Gamma_a(e_b,e_c)+D_aD_b(e_c)\bigr)$.
\end{proposition}
\begin{proof}
Since $E^{(p)}(1)=0$ for $p\ge1$, $[(s-1)^2]E_aE_bE_c(1)=E_a^{(2)}(e_be_c)+D_aD_b(e_c)+e_aE_b^{(2)}(e_c)$, and $E_a^{(2)}(e_be_c)=\Gamma_a(e_b,e_c)+e_bE_a^{(2)}(e_c)+e_cE_a^{(2)}(e_b)$. Each of $e_bE_a^{(2)}(e_c)$, $e_cE_a^{(2)}(e_b)$, $e_aE_b^{(2)}(e_c)$ is a part $e_{\lambda_i}$ times a function of the type of the other two parts (\cref{prop:block}), so every $\mu$ in its support has $\kappa(\lambda,\mu)\ge2$.
\end{proof}

\begin{theorem}[Constant-term adjoint formula]\label{thm:CT}
For $g\in\Lam_a$ and $F\in\Lam$,\[ \varphi_a(t)\,\langle T_ag,F\rangle_t=\mathrm{CT}\bigl[Z\,g(z)\,F(z_1^{-1},\dots,z_a^{-1})\,\Omega(z)\bigr], \]where $Z=z_1\cdots z_a$ and $\Omega(z)=\prod_{1\le i<j\le a}\frac{z_j-z_i}{z_j-tz_i}$, expanded in $z_i/z_j$ ($i<j$).
\end{theorem}
\begin{proof}
Let $G=T_ag$. In exactly $a$ variables only $A=[a]$ survives and $c_{[a]}=1$, so $G(z)=Zg(z)$; and $G\in\mathrm{span}\{P_\mu:\ell(\mu)=a\}$ by \cref{lem:parabolic}. We read $\mathrm{CT}$ of a Laurent polynomial times $\Omega$ either formally or as the integral over the torus $|z_i|=1$ for $|t|<1$; the two agree because the expansion of $\Omega$ converges there, and both sides of the theorem are rational in $t$.

\emph{Step 1.} For $\alpha\in\mathbb{Z}^a_{\ge0}$, $\langle G,q_\alpha\rangle_t=[y^\alpha]G(y_1,\dots,y_a)$, with $q_\alpha=\prod q_{\alpha_i}$: pair $G$ with the Cauchy identity $\sum_\mu Q_\mu(x)P_\mu(y)=\prod_{i,j}\frac{1-tx_iy_j}{1-x_iy_j}=\sum_\alpha q_\alpha(x)y^\alpha$ \cite[III (4.2), (4.4)]{Macdonald95}, using $\langle P_\lambda,Q_\mu\rangle_t=\delta_{\lambda\mu}$ \cite[III (4.9)]{Macdonald95}.

\emph{Step 2.} For $\ell(\lambda)\le a$ (padded), the raising-operator formula \cite[III (2.15$'$)]{Macdonald95} gives $Q_\lambda=\sum_\beta c_\beta q_{\lambda+\beta}$ with $\sum_\beta c_\beta y^{-\beta}=\prod_{i<j}\frac{1-y_j/y_i}{1-ty_j/y_i}$; operators involving positions beyond $\ell(\lambda)$ only produce negative entries, where $q$ and $[y^\cdot]$ both vanish. With step 1 and the relabelling $z_k=y_{a+1-k}$,
\[ \langle G,Q_\lambda\rangle_t=\mathrm{CT}\bigl[G(z)\,z^{-\lambda^{\mathrm{rev}}}\,\Omega(z)\bigr],\qquad\lambda^{\mathrm{rev}}=(\lambda_a,\dots,\lambda_1). \]

\emph{Step 3.} Let $\Delta=\prod_{i\ne j}\frac{1-z_i/z_j}{1-tz_i/z_j}$ and $\theta=\prod_{i<j}\frac{z_i-tz_j}{z_i-z_j}$, so that $\Delta\theta=\Omega$. For symmetric $f$, $f\Omega$ is analytic on the torus and $f\Delta$ is symmetric, so averaging over $S_a$ and using $\sum_ww(\theta)=v_a(t)=[a]_t!$ \cite[III (1.4)]{Macdonald95} gives $\mathrm{CT}[f\Omega]=\frac{v_a}{a!}\mathrm{CT}[f\Delta]$.

\emph{Step 4.} By \cite[III (2.1)--(2.2)]{Macdonald95}, $Q_\lambda(z^{-1})=\frac{b_\lambda}{v_\lambda}\sum_ww\bigl(z^{-\lambda}\tilde\theta\bigr)$ with $\tilde\theta=\prod_{i<j}\frac{z_j-tz_i}{z_j-z_i}$ and $v_\lambda=\prod_{i\ge0}v_{m_i(\lambda)}$, $m_0=a-\ell(\lambda)$. Since $\tilde\theta\Delta=\prod_{i<j}\frac{z_i-z_j}{z_i-tz_j}$, step 3, symmetry, and the relabelling $z_k\leftrightarrow z_{a+1-k}$ give
\[ \mathrm{CT}\bigl[G\,Q_\lambda(z^{-1})\,\Omega\bigr]=v_a\frac{b_\lambda}{v_\lambda}\,\mathrm{CT}\bigl[G\,z^{-\lambda^{\mathrm{rev}}}\Omega\bigr]=v_a\frac{b_\lambda}{v_\lambda}\langle G,Q_\lambda\rangle_t . \]

\emph{Step 5.} Write $F=\sum_\lambda F_\lambda Q_\lambda$. Terms with $\ell(\lambda)>a$ vanish on both sides, and $\langle G,Q_\lambda\rangle_t=0$ unless $\ell(\lambda)=a$, in which case $m_0=0$ and $v_a\,b_\lambda/v_\lambda=v_a(1-t)^a=\varphi_a$. Hence $\mathrm{CT}[GF(z^{-1})\Omega]=\varphi_a\langle G,F\rangle_t$.
\end{proof}

\begin{theorem}[Two-point formula]\label{thm:2pt}
Let $g\in\Lam_a$ be homogeneous of degree $d$, $n=a+d$, $x,y\ge1$, $x+y=n$, and $\Phi_a(g;x,y):=\langle T_ag,p_xp_y\rangle$. For $1\le A\le a-1$ put $B=a-A$ and $G_A(w)=\sum_kG_{A,k}w^k:=g(1,t,\dots,t^{A-1},w,wt,\dots,wt^{B-1})$. Then
\begin{multline*}
\Phi_a(g;x,y)=(-1)^a(1-t^x)(1-t^y)\Biggl\{\sum_{A=1}^{a-1}t^{-AB}\Bigl[\frac{G_{A,y-B}}{(1-t^A)(1-t^B)}\\+\frac{1}{1-t^a}\sum_{m\ge1}(t^{-Am}-t^{Bm})\,G_{A,y-B-m}\Bigr]\\
-\frac{g(1,t,\dots,t^{a-1})}{1-t^a}\sum_{j=0}^{a-1}t^{-jy}\Biggr\}.
\end{multline*}
\end{theorem}
\begin{proof}
Let $S(u):=\sum_{i=1}^a\frac{1}{uz_i-1}=\sum_{x\ge1}u^{-x}p_x(z^{-1})$. By \cref{thm:CT} and homogeneity,
\[ \mathcal R(u,v):=\frac1{\varphi_a}\mathrm{CT}\bigl[g\,Z\,S(u)S(v)\,\Omega\bigr]=\sum_{x+y=n}\frac{\Phi_a(g;x,y)}{(1-t^x)(1-t^y)}u^{-x}v^{-y}=u^{-n}\rho(w),\qquad w=u/v, \]
with $\rho$ a polynomial, so $\Phi_a(g;x,y)=(1-t^x)(1-t^y)[w^y]\rho$.

\emph{String expansion.} Take $|t|$ small, $u,v$ generic, and radii $\max(|u|^{-1},|v|^{-1})<r_1<\dots<r_a$; the constant term is the integral of $g\,S(u)S(v)\,\Omega\prod_mdz_m/(2\pi i)$ over $|z_m|=r_m$. Expand $S(u)S(v)$ into its $a^2$ summands and integrate $z_1,\dots,z_a$ in turn. When $z_m$ is integrated, its poles inside $|z_m|=r_m$ are at $1/u$ or $1/v$ if $z_m$ carries the $u$- or $v$-factor, and at $tz_l$ for $l<m$ (from the factor $\frac{z_m-z_l}{z_m-tz_l}$); the poles $z_j/t$, $j>m$, lie outside, and $g$ is a polynomial. The numerator $\prod_{l<m}(z_m-z_l)$ kills the residue at $tz_l$ if some earlier $z_{l'}$ already equals $tz_l$. Hence the evaluation points form \emph{$t$-strings} $q,qt,\dots,qt^{A-1}$ (at increasing positions), each started by a \emph{seed} at $1/u$ or $1/v$, and position $1$ is always a seed. The surviving configurations are: (i) a word $\omega\in\{U,V\}^a$ containing both letters, the $u$-factor at the first $U$ and the $v$-factor at the first $V$, the $(\alpha+1)$-st $U$ at $t^\alpha/u$ and the $(\beta+1)$-st $V$ at $t^\beta/v$; (ii) the word $U^a$, all points $t^{m-1}/u$, with the $v$-factor at any position $j+1$, where it is simply evaluated; (iii) the mirror image $V^a$.

\emph{String weight.} A string of length $A$ with base $q$ contributes $(-1)^{A-1}\varphi_{A-1}(t)\,q^A$: the seed residue is $q$; a successor pair contributes the residue $qt^\alpha(t-1)$; a non-successor pair at string distance $\delta\ge2$ contributes $\frac{t^\delta-1}{t^\delta-t}$; and $\prod_{\delta=2}^{A-1}\bigl(\frac{1-t^\delta}{t(1-t^{\delta-1})}\bigr)^{A-\delta}=t^{-\binom{A-1}2}\varphi_{A-1}/(1-t)^{A-1}$.

\emph{Interaction.} For a $U$-string of length $A$ and a $V$-string of length $B=a-A$, the cross factors are $\frac{wt^{\beta-\alpha}-1}{wt^{\beta-\alpha}-t}$ ($U_\alpha$ before $V_\beta$) and $\frac{1-wt^{\beta-\alpha}}{1-wt^{\beta-\alpha+1}}$ ($V_\beta$ before $U_\alpha$), and their sum over shuffles is the function $\mathrm{Sh}_{A,B}(w)$ of \cref{lem:shuffle}.

\emph{Assembling.} Since $g(z^*)=u^{-d}G_A(w)$ for configuration (i) and $u^{-d}g(\pi_a)$ for (ii), where $\pi_a=(1,t,\dots,t^{a-1})$, and $v^{-B}=u^{-B}w^B$,
\[ \varphi_a\rho(w)=\sum_{A=1}^{a-1}(-1)^a\varphi_{A-1}\varphi_{B-1}\,w^B\,\mathrm{Sh}_{A,B}(w)\,G_A(w)+(-1)^{a-1}\varphi_{a-1}\,g(\pi_a)\sum_{j=0}^{a-1}\frac{w}{t^j-w}+O(w^n), \]
the last term coming from (iii). This is an identity of rational functions of $w$, each term analytic at $w=0$, so $[w^y]$ may be taken termwise. By \cref{lem:shuffle}, $\mathrm{Sh}_{A,B}(w)=\qbin aA t^{-AB}\bigl(1+\kappa_{AB}\sum_{m\ge1}(t^{-Am}-t^{Bm})w^m\bigr)$ with $\kappa_{AB}=\frac{(1-t^A)(1-t^B)}{1-t^a}$, and $\qbin aA\varphi_{A-1}\varphi_{B-1}=\frac{\varphi_a}{(1-t^A)(1-t^B)}$, $\frac{\varphi_{a-1}}{\varphi_a}=\frac1{1-t^a}$. For $1\le y\le n-1$ this gives the formula.
\end{proof}

\begin{lemma}[Shuffle identity]\label{lem:shuffle}
\[ \mathrm{Sh}_{A,B}(w)=\qbin{A+B}{A}\,t^{-AB}\,\frac{(1-w)(1-t^{B-A}w)}{(1-t^{-A}w)(1-t^Bw)}, \]\[ \frac{(1-w)(1-t^{B-A}w)}{(1-t^{-A}w)(1-t^Bw)}=1+\kappa_{AB}\sum_{m\ge1}(t^{-Am}-t^{Bm})w^m. \]
\end{lemma}
\begin{proof}
Induction on $A+B$; for $A=0$ or $B=0$ both sides are $1$. Condition on the last letter. If it is $U_{A-1}$, its cross factors telescope to $\frac{1-wt^{1-A}}{1-wt^{B+1-A}}$; if it is $V_{B-1}$, to $t^{-A}\frac{1-wt^{B-1}}{1-wt^{B-1-A}}$. So $\mathrm{Sh}_{A,B}=\frac{1-wt^{1-A}}{1-wt^{B+1-A}}\mathrm{Sh}_{A-1,B}+t^{-A}\frac{1-wt^{B-1}}{1-wt^{B-1-A}}\mathrm{Sh}_{A,B-1}$. With $X=t^A$, $Y=t^B$, $\qbin{A+B-1}{A-1}/\qbin{A+B}{A}=\frac{1-X}{1-XY}$ and $\qbin{A+B-1}{A}/\qbin{A+B}A=\frac{1-Y}{1-XY}$, the claimed closed form satisfies this recursion if and only if $Y(1-X)(1-w/X)+(1-Y)(1-Yw)=(1-XY)(1-wY/X)$, and both sides equal $1-XY-wY/X+wY^2$. For the expansion: the left side is $1$ at $w=0$, tends to $1$ at $\infty$, and has simple poles with residue $-t^A\kappa_{AB}$ at $w=t^A$ and $t^{-B}\kappa_{AB}$ at $w=t^{-B}$; so does the right side.
\end{proof}
\begin{remark}[Relation to Green polynomials]
Write $p_\mu=\sum_\lambda X^\lambda_\mu(t)P_\lambda(x;t)$. For $g=P_\rho$ and $\lambda=\rho+1^a$ one has $\Phi_a(P_\rho;x,y)=(1-t^x)(1-t^y)X^{\lambda}_{(x,y)}(t)/b_\lambda(t)$, so \cref{thm:2pt} evaluates Green polynomials at two-part classes; the two $t$-strings are the Frobenius orbits of a maximal torus of type $(x,y)$ \cite{Green55}. \Cref{thm:CT} is Jing's constant-term expression for the matrix element $\langle Q_\lambda,p_\mu\rangle$ \cite{Jing91}, cf.\ \cite[(2.17), (2.20)]{JingLiu21}. Jing and Liu \cite[Thm.~2.7]{JingLiu21} expand the same matrix element as a nested sum valid for all $\mu$, extracting one variable at a time; at two-part $\mu$ that sum does not close up, since its inner sums run over classes of arbitrary length. Their closed forms \cite[Thm.~2.10, (2.37)--(2.40)]{JingLiu21} restrict the Hall--Littlewood index $\lambda$, not the class $\mu$; their Murnaghan--Nakayama rule \cite[Thm.~3.2]{JingLiu21} generalises a formula of Morris \cite{Morris77}\footnote{We have not seen \cite{Morris77} first-hand; its scope is as described in \cite[p.~12, before Thm.~3.2]{JingLiu21}.}, its case $\ell(\lambda)=2$. We evaluate by residues instead, and at $\ell(\mu)=2$ the poles organise into two $t$-strings. The formula is explicit, with no sum over classes, and linear in the two-string specialisation $G_A$ of $g$; for general $\rho$, $G_A$ is itself a finite (branching) sum, so \cref{thm:2pt} is not a product formula.
\end{remark}
\begin{example}[Two rows]\label{ex:tworow}
For $a=2$, only $A=B=1$ occurs, and $G_1(w)=P_\rho(1,w;t)=w^{\rho_2}\frac{(1-tw)-w^r(w-t)}{1-w}$ with $r=\rho_1-\rho_2\ge1$ ($G_1=w^{\rho_2}$ if $r=0$). Substituting into \cref{thm:2pt} gives, for $\lambda=(\lambda_1,\lambda_2)$, $y\le x$, and $m_{xy}=2$ if $x=y$, $1$ otherwise,
\[ X^{(\lambda_1,\lambda_2)}_{(x,y)}(t)=\begin{cases}(t-1)\,t^{\lambda_2-1-y}(1+t^{y}), & y<\lambda_2,\\ m_{xy}-(1-t)\,t^{\lambda_2-1}, & y=\lambda_2,\\ (t-1)\,t^{\lambda_2-1}, & y>\lambda_2,\end{cases} \]
in agreement with a Kostka--Foulkes computation for $|\lambda|\le10$. This is not new: for $\lambda=(n-k,k)$ and every class $\rho$, Jing and Liu \cite[\S2, display after (2.32)]{JingLiu21} give $X^\lambda_\rho=\pi_k+(t-1)\sum_{j<k}t^{k-1-j}\pi_j$, with $\pi_j$ the number of $j$-subsets fixed by $\rho$ (their $D^{(j)}(\rho)$). A short proof: $X^\lambda_\rho=\sum_\nu\chi^\nu_\rho K_{\nu\lambda}(t)$ \cite[III (7.6$'$)]{Macdonald95}, where $K_{(n-j,j),\lambda}=t^{k-j}$ (one tableau, and $K$ is monic of degree $n(\lambda)-n(\nu)$ \cite[III (6.5)]{Macdonald95}) and $\chi^{(n-j,j)}_\rho=\pi_j-\pi_{j-1}$ by Young's rule. On the diagonal class $(x,y)=\lambda$ it gives $t^{\lambda_2}-t^{\lambda_2-1}+m_{xy}$, which for $\lambda_1>\lambda_2$ with $\lambda_2\ge3$ odd equals $-1$ at $t=-1$ and $1$ at $t=0$. A root in $(-1,0)$ rules out a product of cyclotomic factors and powers of $t$, and Green polynomials at two-part classes admit no product formula even for two rows; for products of the form $t^c\prod_i(1-t^{d_i})$ this was observed independently in~\cite[Thm.~D]{Clio26}.
\end{example}

\begin{theorem}[Closed $\ell=3$, $\kappa=1$ leads]\label{thm:v2}
Let $\lambda=(a,b,c)$ be any ordering of a length-$3$ partition, $\mu=(x,y)$ with $\kappa(\lambda,\mu)=1$, $n=x+y$, $m_{xy}$ as in \cref{ex:tworow}. Put \[\Xi_a(r,q):=\lin T_a(p_rp_q)=(-1)^{r+q}\frac{[a+r+q]}{[a]}[a]_{t^r}[a]_{t^q}\] and $U_a(b,c;x,y):=\bigl((-1)^n\Phi_a(p_bp_c;x,y)-\Xi_a(b,c)\bigr)/m_{xy}$. Then
\begin{multline*}
[(s-1)^2]\,c_{\lambda\mu}=(-1)^{b+c}U_a(b,c;x,y)+\sum_{\substack{1\le r<b\\ \{b-r,\,a+r+c\}=\{x,y\}}}(-1)^{r+c}\Xi_a(r,c)\\
+\sum_{\substack{1\le q<c\\ \{c-q,\,a+b+q\}=\{x,y\}}}(-1)^{b+q}\Xi_a(b,q)+[e_xe_y]\,D_a(M_{bc}),
\end{multline*}
with $D_a$ the derivation $D_a(e_j)=M_{aj}$ of \cref{prop:M}.
\end{theorem}
\begin{lemma}\label{lem:pair2}
For $G$ homogeneous of degree $n=x+y$, $\langle G,p_xp_y\rangle=(-1)^n\bigl(m_{xy}[e_xe_y]G+\lin G\bigr)$.
\end{lemma}
\begin{proof}
$\langle e_\nu,p_xp_y\rangle=\langle\Delta e_\nu,p_x\otimes p_y\rangle$ since $p_x,p_y$ are primitive, and $\langle e_m,p_m\rangle=(-1)^{m-1}$. Each tensor factor must be a single $e$, so $\ell(\nu)\ge3$ gives $0$, $\langle e_n,p_xp_y\rangle=(-1)^n$, $\langle e_xe_y,p_xp_y\rangle=(-1)^nm_{xy}$, and other two-part $\nu$ give $0$.
\end{proof}
\begin{proof}[Proof of \cref{thm:v2}]
By \cref{prop:reduce}, $[(s-1)^2]c_{\lambda\mu}=[e_xe_y]\bigl(\Gamma_a(e_b,e_c)+D_a(M_{bc})\bigr)$. In $\Gamma_a(e_b,e_c)=\sum_{r,q}(-1)^{r+q}e_{b-r}e_{c-q}T_a(p_rp_q)$, terms with $r<b$ and $q<c$ have at least three factors in $I$ and do not contribute. For $r<b$, $q=c$, only the single-$e$ part $\Xi_a(r,c)e_{a+r+c}$ of $T_a(p_rp_c)$ (\cref{lem:linT}) contributes; symmetrically for $r=b$, $q<c$. For $r=b$, $q=c$ the contribution is $(-1)^{b+c}[e_xe_y]T_a(p_bp_c)=(-1)^{b+c}U_a(b,c;x,y)$ by \cref{lem:pair2} and $\lin T_a(p_bp_c)=\Xi_a(b,c)$.
\end{proof}
\begin{corollary}\label{cor:v2}
Every lead at valuation $2$ ($\kappa=\ell-2$, any $\ell$) has a closed formula: by \cref{thm:blockmult} it is a sum of products of pair leads (\cref{prop:M}) and $\ell=3$, $\kappa=1$ leads (\cref{thm:G}, \cref{cor:column}, \cref{thm:v2}).
\end{corollary}
\begin{example}\label{ex:v2}
$(3,3,3)\to(7,2)$: $\Lead=2t^{13}+3t^{12}+3t^{11}+6t^{10}+6t^9+6t^8+9t^7+6t^6+3t^5+9t^4+6t^3+3t+4$ (agrees with direct computation). Leads are \emph{not} positive in general: $(4,4,2)\to(7,3)$ has $\Lead=(t+1)(t^2+1)(2t^8+t^7+t^6+3t^4+t^3+t^2-t+2)$.
\end{example}
\paragraph{Computer checks.} All statements in this section were checked against an independent implementation of Hikita's subset formula: \cref{thm:v2} on $16$ pairs with $n\le12$ (both examples among them) and on $27$ pairs in all $3$ orderings with $n\le10$; \cref{thm:2pt} in $229$ cases ($a\le5$); and the Green-polynomial dictionary symbolically for $|\lambda|\le8$ and numerically for $|\lambda|\le9$.


%% ===================================================================
\section{Open problems}\label{sec:open}
%% ===================================================================
\begin{enumerate}
  \item \textbf{Valuation $\ge3$.} Connected leads with $\ell\ge4$, $\kappa=1$. Hierarchy: $v=1\leftrightarrow p_n\leftrightarrow$ one $t$-string (\cref{lem:linT}); $v=2\leftrightarrow p_xp_y\leftrightarrow$ two strings (\cref{thm:2pt}). Does the three-string shuffle sum factor pairwise?
  \item \textbf{Zero set at generic $t$.} At $t=-2$ seven pairs with $n=6$ have $\val>\ell-\kappa$; for coarsenings this is cancellation between histories, e.g.\ $\Lead_{(1,1,1),(3)}=[3]_t(2+t)$. Describe the zero set of $[(s-1)^{\ell-\kappa}]c_{\lambda\mu}$ in $t$.
  \item \textbf{Counting.} A formula for $N(\lambda,\mu)$ in \cref{thm:blocklaw}; for $\ell(\lambda)=3$ and $\mu=(n-1,1)$ with $\kappa=1$ we observe $\Lead_{\lambda\mu}(1)=2n^2-6n+3$, independent of $\lambda$ ($n\le10$; computed only).
  \item \textbf{Positivity.} Leads are not $t$-positive (\cref{ex:v2}), so no ``$t$-count of minimal flows'' can hold. Is there a signed or cancellation-free model?
  \item \textbf{Cumulant link.} Is there a Macdonald-cumulant construction whose leads are the $\star$-leads for all $t$, i.e.\ cumulants of $e_k\mapsto t^{\binom k2}$?
\end{enumerate}

\printbibliography

\end{document}
